% Dokumentklasse
\documentclass[11pt,titlepage]{report} % Für längere Arbeiten mit Kapiteln und Abschnitten, geeignet für Berichte und Dokumentationen.

% Sprachunterstützung
\usepackage[ngerman]{babel} % Sorgt für sprachspezifische Regeln wie Silbentrennung und deutsche Bezeichnungen.

% Seitenlayout
\usepackage[a4paper,lmargin={3.5cm},rmargin={2.5cm},tmargin={3cm},bmargin={2cm},headsep=0.5cm,headheight=1.5cm]{geometry} % Einstellung von Seitenrändern für A4-Papier.

% Mathematische Pakete
\usepackage{amssymb} % Enthält zusätzliche mathematische Symbole.
\usepackage{amsthm} % Unterstützung für Theoreme, Definitionen und Beweise.
\usepackage{amsmath} % Verbesserte Darstellung von mathematischen Formeln.

% Grafiken und Bilder
\usepackage{graphicx} % Ermöglicht das Einfügen von Bildern und Grafiken.
\usepackage{svg} % Ermöglicht das Einfügen von SVG-Grafiken.

% Typografie
\usepackage{microtype} % Verbessert die Lesbarkeit und Silbentrennung.
\usepackage{fontspec} % Für die Verwendung von TrueType- und OpenType-Schriften.

% Kopf- und Fußzeilen
\usepackage{fancyhdr} % Ermöglicht die Gestaltung von Kopf- und Fußzeilen.

% Farben
\usepackage{xcolor} % Definiert und verwendet Farben im Dokument.

% Weitere mathematische Symbole
\usepackage{amsfonts} % Zusätzliche Schriftarten für mathematische Symbole.

% Tabellen
\usepackage{booktabs} % Verbesserte Gestaltung von Tabellen.
\usepackage{multirow} % Ermöglicht Zellen über mehrere Zeilen in Tabellen.
\usepackage{enumitem} % Anpassung von Listen und Aufzählungen.
\usepackage{tabularx} % Erweiterte Tabellenfunktionen, insbesondere für Spaltenbreiten.

% Literaturverwaltung
\usepackage[notes,backend=biber,autocite=footnote,ibidtracker=false,dashed=false]{biblatex-chicago} % Für Literaturverzeichnisse nach Chicago-Stil (Notes and Bibliography).

% Zitate
\usepackage{csquotes} % Korrekte Darstellung von Zitaten, z. B. Anführungszeichen.

% Fußnoten-Formatierung
\usepackage[hang,flushmargin,bottom,multiple]{footmisc} % Keine Einrückung, mehr Abstand zwischen Fußnoten
\setlength{\footnotesep}{12pt} % Abstand zwischen Fußnoten

% Zeilenabstand
\usepackage[onehalfspacing]{setspace} % Erzeugt 1,5-fachen Zeilenabstand.

% Hyperlinks
\usepackage[hidelinks]{hyperref}% Fügt klickbare Hyperlinks ein, ohne sie farblich hervorzuheben.

% Positionierung
\usepackage{float} % Präzise Platzierung von Tabellen und Abbildungen.

% Glossare und Abkürzungen
\usepackage[toc=false]{glossaries-extra} % Verwaltung von Glossaren ohne Inhaltsverzeichnis-Eintrag.

% Tabellenfarben
\usepackage{colortbl} % Ermöglicht das Einfärben von Tabellenzellen.

% Tabellen- und Spaltenformatierung
\usepackage{array} % Zusätzliche Optionen für Tabellen und Spalten.

% Diagramme und Visualisierungen
\usepackage{tikz} % Erstellung von Vektorgraphiken und Diagrammen.
\usepackage{pgfplots} % Erstellung von Plots und Grafiken.

\usetikzlibrary{shapes,arrows,positioning,fit,backgrounds,calc,matrix,decorations.pathreplacing,decorations.markings,shapes.geometric,mindmap}

% PDF-Integration
\usepackage{pdfpages} % Einfügen von ganzen PDF-Dokumenten.

% Querformat
\usepackage{pdflscape} % Ermöglicht die Nutzung von Seiten im Querformat.

% Nummerierungsanpassung
\usepackage{chngcntr} % Ermöglicht Änderungen der Nummerierung von Kapiteln und Abschnitten.
\counterwithout{figure}{chapter} % Abbildungen durchgehend nummerieren (ohne Kapitelpräfix)
\counterwithout{footnote}{chapter} % Fußnoten durchgehend nummerieren (ohne Kapitelpräfix)
\counterwithout{table}{chapter} % Tabellen durchgehend nummerieren (ohne Kapitelpräfix)

% Inhaltsverzeichnis
\usepackage{tocloft} % Anpassung des Layouts des Inhaltsverzeichnisses.

% Titelgestaltung
\usepackage{titlesec} % Anpassung von Kapitel- und Abschnittstiteln.

% Ausrichtung
\usepackage{ragged2e} % Ermöglicht den Blocksatz oder Flattersatz.

% Coding
\usepackage{listings} % Für die Darstellung von Quellcode.

\pgfplotsset{compat=1.18}

% Pfade für Bilder und andere Ressourcen
\graphicspath{{bilder/}}

% Arial als Hauptschrift (11pt)
\setmainfont{Arial}

% Bibliographie
\addbibresource{quellen.bib}

% Farben
\definecolor{link}{HTML}{0000EE}

% Absatzformatierung
\setlength{\parindent}{0pt}           % Keine Einrückung am Absatzbeginn
\setlength{\parskip}{6pt}             % 6pt Abstand nach jedem Absatz

% Anpassung \listoffigures
\renewcommand{\cftloftitlefont}{\normalsize\bfseries}
\renewcommand{\cftbeforeloftitleskip}{10pt}
\renewcommand{\cftafterloftitleskip}{10pt}

% Anpassung \listoftables
\renewcommand{\cftlottitlefont}{\normalsize\bfseries}
\renewcommand{\cftbeforelottitleskip}{10pt}
\renewcommand{\cftafterlottitleskip}{10pt}

% Anpassung \tableofcontents
\renewcommand{\cfttoctitlefont}{\normalsize\bfseries}
\renewcommand{\cftbeforetoctitleskip}{10pt}
\renewcommand{\cftaftertoctitleskip}{10pt}

% Füge Punkte für Kapitel im Inhaltsverzeichnis hinzu
\renewcommand{\cftchapaftersnum}{.}
\renewcommand{\cftchapdotsep}{\cftdotsep}
\renewcommand{\cftsecaftersnum}{.}

% Präfixe für Abbildungs- und Tabellenverzeichnis
\renewcommand{\cftfigpresnum}{Abb.~}
\renewcommand{\cfttabpresnum}{Tab.~}

% Abstand nach Nummern im Abbildungs- und Tabellenverzeichnis
\setlength{\cftfignumwidth}{4em}
\setlength{\cfttabnumwidth}{4em}

% Überschriftenformatierung
\titleformat{\chapter}
  {\fontsize{11}{16.5}\bfseries\raggedright}
  {\thechapter.}{1em}{}
\titlespacing*{\chapter}{0pt}{0pt}{6pt}

\titleformat{\section}
  {\fontsize{11}{16.5}\bfseries\raggedright}
  {\thesection.}{1em}{}
\titlespacing*{\section}{0pt}{0pt}{6pt}

\titleformat{\subsection}
  {\fontsize{11}{16.5}\bfseries\raggedright}
  {\thesubsection.}{1em}{}
\titlespacing*{\subsection}{0pt}{0pt}{6pt}

\titleformat{\subsubsection}
  {\fontsize{11}{16.5}\bfseries\raggedright}
  {\thesubsubsection.}{1em}{}
\titlespacing*{\subsubsection}{0pt}{0pt}{6pt}

% Blocksatz für Fließtext
\justifying

% Bibliographie-Einstellungen für Chicago-Stil
\DeclareNameAlias{sortname}{family-given}
\DefineBibliographyStrings{german}{andothers = {{et\,al\adddot}},}

% Benutzerdefinierte Zitatbefehle
\newcommand{\customcite}[2][]{%
  \ifblank{#1}{%
    \autocite{#2}%
  }{%
    \autocite[#1]{#2}%
  }%
}

% Fix lange URLs
\newcommand*\oldurlbreaks{}
\let\oldurlbreaks=\UrlBreaks
\renewcommand{\UrlBreaks}{\oldurlbreaks
  \do\a\do\b\do\c\do\d\do\e\do\f\do\g\do\h\do\i\do\j\do\k\do\l\do\m\do\n\do\o\do\p\do\q\do\r\do\s\do\t\do\u\do\v
  %\do\w%
  \do\x\do\y\do\z
  \do?\do&\do_
  \do\0\do\1\do\2\do\3\do\4\do\5\do\6\do\7\do\8\do\9
  \do:\do/
  }
\urlstyle{same}

% lstlisting Einstellungen 
\lstset{  
  tabsize=2, 
  showspaces=false, 
  showstringspaces=false, 
  float=[htb], 
  captionpos=b, 
  basicstyle=\footnotesize, 
  numberblanklines=false,
  numberbychapter=false
} 

% Befehle für Links und Verweise
\newcommand{\sct}[2]{\textcolor{link}{\hyperref[#1]{#2}}}
\newcommand{\gl}[1]{\gls{#1}}

% Promptverzeichnis erstellen
\newcommand{\listpromptname}{Promptverzeichnis}
\newlistof{prompt}{prm}{\listpromptname}
% Formatierung des Promptverzeichnis-Titels (wie bei anderen Verzeichnissen)
\renewcommand{\cftprmtitlefont}{\normalsize\bfseries}
\renewcommand{\cftbeforeprmtitleskip}{10pt}
\renewcommand{\cftafterprmtitleskip}{10pt}

% Befehl für Prompt-Referenzen
\newcounter{promptcounter}
\newcommand{\promptentry}[1]{%
  \refstepcounter{promptcounter}%
  \addcontentsline{prm}{prompt}{\protect\numberline{\thepromptcounter}#1}%
  #1%
}

% Inhaltsverzeichnis-Einrückung
\setlength{\cftsubsecindent}{1.5em}
\setlength{\cftsubsecindent}{3.0em}

% Anhangsverzeichnis erstellen
\newcommand{\listappendixname}{Verzeichnis der Anhänge}
\newlistof{appendix}{app}{\listappendixname}

% Formatierung des Anhangsverzeichnis-Titels (wie bei anderen Verzeichnissen)
\renewcommand{\cftapptitlefont}{\normalsize\bfseries}
\renewcommand{\cftbeforeapptitleskip}{10pt}
\renewcommand{\cftafterapptitleskip}{10pt}

% Setze die Breite für Anhangsnummern im Verzeichnis
\newlength{\appnumwidth}
\setlength{\appnumwidth}{6em}

\newcommand{\appendixentry}[1]{%
  \refstepcounter{appendix}%
  \addcontentsline{app}{appendix}{\protect\makebox[\appnumwidth][l]{Anhang \Alph{appendix}:}#1}%
}

% Glossar und Abkürzungen
\makenoidxglossaries
\newacronym{ack}{ACK}{Acknowledgment}
\newacronym{api}{API}{Application Programming Interface}
\newacronym{aws}{AWS}{Amazon Web Services}
\newacronym{cd}{CD}{Continuous Delivery}
\newacronym{cfs}{CFS}{Completely Fair Scheduler}
\newacronym{ci}{CI}{Continuous Integration}
\newacronym{cncf}{CNCF}{Cloud Native Computing Foundation}
\newacronym{cni}{CNI}{Container Network Interface}
\newacronym{cpu}{CPU}{Central Processing Unit}
\newacronym{crd}{CRD}{Custom Resource Definition}
\newacronym{csv}{CSV}{Comma-Separated Values}
\newacronym{dns}{DNS}{Domain Name System}
\newacronym{dram}{DRAM}{Dynamic Random-Access Memory}
\newacronym{ec2}{EC2}{Elastic Compute Cloud}
\newacronym{grpc}{gRPC}{Google Remote Procedure Call}
\newacronym{http}{HTTP}{Hypertext Transfer Protocol}
\newacronym{https}{HTTPS}{Hypertext Transfer Protocol Secure}
\newacronym{iac}{IaC}{Infrastructure as Code}
\newacronym{icmp}{ICMP}{Internet Control Message Protocol}
\newacronym{imds}{IMDSv2}{Instance Metadata Service Version 2}
\newacronym{iot}{IoT}{Internet of Things}
\newacronym{json}{JSON}{JavaScript Object Notation}
\newacronym{mttr}{MTTR}{Mean Time to Recovery}
\newacronym{oom}{OOM}{Out of Memory}
\newacronym{promql}{PromQL}{Prometheus Query Language}
\newacronym{ram}{RAM}{Random-Access Memory}
\newacronym{rest}{REST}{Representational State Transfer}
\newacronym{rfc}{RFC}{Request for Comments}
\newacronym{rst}{RST}{Reset (TCP-Steuerflag)}
\newacronym{rto}{RTO}{Retransmission Timeout}
\newacronym{scp}{SCP}{Secure Copy Protocol}
\newacronym{sha}{SHA}{Secure Hash Algorithm}
\newacronym{sre}{SRE}{Site Reliability Engineering}
\newacronym{ssd}{SSD}{Solid-State Drive}
\newacronym{ssh}{SSH}{Secure Shell}
\newacronym{ssot}{SSOT}{Single Source of Truth}
\newacronym{sut}{SUT}{System Under Test}
\newacronym{tcp}{TCP}{Transmission Control Protocol}
\newacronym{tls}{TLS}{Transport Layer Security}
\newacronym{tsdb}{TSDB}{Time Series Database}
\newacronym{ui}{UI}{User Interface}
\newacronym{vcpu}{vCPU}{Virtual Central Processing Unit}
\newacronym{vpc}{VPC}{Virtual Private Cloud}
\newacronym{wan}{WAN}{Wide Area Network}
\newacronym{yaml}{YAML}{YAML Ain't Markup Language}

% Trenner ändern - Bindestrich zwischen Abkürzung und Erklärung
\renewcommand*{\glsnamefont}[1]{#1 --}

% Anpassung der Kopfzeilen
% Behält die Nummer bei nummerierten Kapiteln, entfernt nur das Wort "Kapitel"
\renewcommand{\chaptermark}[1]{\markboth{\thechapter. #1}{}}

% Seitenstil für plain Seiten
\fancypagestyle{plain}{%
  \fancyhf{} % Alle Kopf- und Fußzeilenfelder löschen
  \lhead{\fontsize{9}{13.5}\selectfont\nouppercase{\leftmark}} % Kapitelüberschrift oben links
  \rhead{\fontsize{9}{13.5}\selectfont\thepage} % Seitenzahl oben rechts
}

% Seitenstil für alle Seiten
\pagestyle{fancy}
  \fancyhf{} % Alle Kopf- und Fußzeilenfelder löschen
  \lhead{\fontsize{9}{13.5}\selectfont\nouppercase{\leftmark}} % Kapitelüberschrift oben links
  \rhead{\fontsize{9}{13.5}\selectfont\thepage} % Seitenzahl oben rechts

% Anpassung der Verzeichnis-Überschriften
\renewcommand{\contentsname}{Inhalt}
\renewcommand{\listfigurename}{Abbildungsverzeichnis}
\renewcommand{\listtablename}{Tabellenverzeichnis}
\renewcommand{\bibname}{Literaturverzeichnis}

% Anpassung der Bezeichnungen für Abbildungen und Tabellen
\addto\captionsngerman{%
  \renewcommand{\figurename}{Abb.}%
  \renewcommand{\tablename}{Tab.}%
}

% Listings für Code
\input{zusatz/listing-definitions}

% Setze PDF-Metadaten
\hypersetup{
  pdftitle={Empirische Untersuchung der Resilienz und Konvergenzlatenz deklarativer GitOps-Control-Planes unter transienten Netzwerkstörungen},
  pdfauthor={Mustermensch},
  pdfsubject={Bachelor-Thesis Informatik (B.Sc.)},
}

\begin{document}

% Leitfaden Titelseite (ohne Kopfzeile, nur für das Leitfaden-Dokument selbst)
%\thispagestyle{empty}
%\begin{titlepage}
    \thispagestyle{empty}
    \hspace{-1cm}\includegraphics[width=0.55\textwidth]{rh-logo}\\[1.9cm]
    
    \vspace{2cm}
    \begin{minipage}{\textwidth}
    \hspace{2.5cm}\begin{minipage}{\dimexpr\textwidth-2.5cm\relax}
    {\Large\textbf{Leitfaden zum Anfertigen \\ wissenschaftlicher Arbeiten}}\\

    
    {\normalsize\fontsize{11pt}{16.5pt}\selectfont Version vom 05. Mai 2025}\\

    \vspace{2cm}
    \textcolor{red}{
        +++ \\
        Löschen Sie diese Seite einfach und verwenden Sie dann die \\
        Vorlagen auf den beiden folgenden Seiten für Ihre Arbeit. \\
        +++ \\
    }

    \end{minipage}
    \end{minipage}
    
    \vfill
    
\end{titlepage}

% Schmutztitel (Vorderdeckel)
\begin{titlepage}
    \thispagestyle{empty}
    \hspace{-1cm}\includegraphics[width=0.55\textwidth]{rh-logo}\\[1.9cm]
    
    \vspace{2cm}
    \begin{minipage}{\textwidth}
    \hspace{2.5cm}\begin{minipage}{\dimexpr\textwidth-2.5cm\relax}
    {\Large\textbf{Fachbereich Informatik und Digitalisierung}}\\
    
    {\Large Studiengang Informatik (B.Sc.)}\\
    \\
    \\
    
    {\Large\textbf{Bachelor-Thesis}}
    
    \parbox{\textwidth}{\Large\raggedright Empirische Untersuchung der Resilienz und Konvergenzlatenz deklarativer GitOps-Control-Planes unter transienten Netzwerkstörungen}\\
    \\
    
    {\normalsize\fontsize{11pt}{16.5pt}\selectfont vorgelegt von}
    
    {\Large Mustermensch}
    
    {\normalsize\fontsize{11pt}{16.5pt}\selectfont Matrikelnr. 111111111}\\
    \\
    
    {\normalsize Wintersemester 2026/27}
    \end{minipage}
    \end{minipage}
    
    \vfill
    
\end{titlepage}


% Leeres weißes Blatt
\newpage
\thispagestyle{empty}
\quad 
\newpage

% Haupttitel
\begin{titlepage}
    \thispagestyle{empty}
    \hspace{-1cm}\includegraphics[width=0.55\textwidth]{rh-logo}\\[1.9cm]
    
    \vspace{2cm}
    \begin{minipage}{\textwidth}
    \hspace{2.5cm}\begin{minipage}{\dimexpr\textwidth-2.5cm\relax}
    {\Large\textbf{Fachbereich Informatik und Digitalisierung}}\\
    
    {\Large Studiengang Informatik (B.Sc.)}\\
    \\
    \\
    
    {\Large\textbf{Bachelor-Thesis}}
    
    \parbox{\textwidth}{\Large\raggedright Empirische Untersuchung der Resilienz und Konvergenzlatenz deklarativer GitOps-Control-Planes unter transienten Netzwerkstörungen}\\
    \\
    
    {\normalsize\fontsize{11pt}{16.5pt}\selectfont vorgelegt von}
    
    {\Large Mustermensch}
    
    {\normalsize\fontsize{11pt}{16.5pt}\selectfont Matrikelnr. 111111111}\\


    {\normalsize\fontsize{11pt}{16.5pt}\selectfont Erstprüfer: Prof. Dr. Mustermensch}

    {\normalsize\fontsize{11pt}{16.5pt}\selectfont Zweitprüfer: Prof. Dr. Mustermensch}\\
    \\
    
    {\normalsize Wintersemester 2026/27}
    \end{minipage}
    \end{minipage}
    
    \vfill
    
\end{titlepage}


% Deaktiviere Einrückungen
\setlength\parindent{0pt}

\begin{onehalfspace}
    % Römische Seitennummerierung für Verzeichnisse
    \pagenumbering{roman}
    
    % Zusammenfassung (deutsch) und Abstract (englisch)
    \chapter*{Zusammenfassung}
\addcontentsline{toc}{chapter}{Zusammenfassung}
\label{sec:zusammenfassung}

Die vorliegende Arbeit untersucht die Resilienz, das Konvergenzverhalten sowie die Ressourcenauslastung deklarativer GitOps-Control-Planes unter transienten Netzwerkstörungen in ressourcenbeschränkten Cloud- und Edge-Umgebungen. Im Fokus steht die Frage, wie sich ein massiver partieller Paketverlust auf der Verbindung zum externen Git-Repository auf die Reconcile-Latenz, die Zustandsübergänge, die Erholungszeit (Mean Time to Recovery, MTTR) sowie den Speicher- und Prozessorbedarf des ArgoCD Application Controllers in einem leichtgewichtigen K3s-Cluster auswirkt. Zur empirischen Beantwortung dieser Fragestellung wurde ein vollständig als Code definiertes, ephemeres Chaos-Engineering-Messframework auf Basis von Terraform und Amazon EC2 konzipiert. Mithilfe deterministischer Fehlerinjektion auf Linux-Kernel-Ebene und hochauflösender Prometheus-Zeitreihenerfassung wird nachgewiesen, dass die Control Plane trotz signifikanter Latenzspitzen während der Störungsphase ihre Selbstheilungsfähigkeit wahrt und ohne Auslagerungsprozesse auf den Sekundärspeicher (Swap-Reinheit) innerhalb der definierten Speichergrenzen stabil konvergiert.

\vspace{0.5cm}

\begin{center}
\textbf{Abstract}
\end{center}

This thesis investigates the resilience, convergence behavior, and resource utilization of declarative GitOps control planes under transient network faults in resource-constrained cloud and edge environments. Specifically, it examines how severe packet loss on the connection to the external Git repository affects reconciliation latency, state transitions, Mean Time to Recovery (MTTR), and the memory and CPU consumption of the ArgoCD Application Controller within a lightweight K3s cluster. To empirically address this research question, a fully automated, ephemeral chaos-engineering benchmarking framework based on Terraform and Amazon EC2 was developed. Using deterministic Linux kernel-level fault injection and high-resolution Prometheus time-series monitoring, the study demonstrates that despite significant latency spikes during the fault window, the control plane preserves its self-healing capability and converges reliably within defined memory boundaries while maintaining complete swap purity.

    \newpage
    
    % Inhaltsverzeichnis
    \phantomsection
    \tableofcontents
    \markboth{}{}
    \newpage

    % Optionale Vorbemerkung
    % Auskommentieren, wenn nicht benötigt
    %\chapter*{Vorbemerkung}
\addcontentsline{toc}{chapter}{Vorbemerkung}
\label{chap:vorbemerkung}

In dieser Arbeit wird das generische Maskulinum verwendet, wenn im Mittelpunkt der Aussagen keine spezifischen Personen stehen, sondern die Funktion, die Personen ausüben. Sämtliche biologischen Geschlechter und subjektiv empfundenen Zuordnungen sind damit gleichermaßen gemeint.

    %\newpage

    % Hauptteil mit arabischer Seitennummerierung
    \pagenumbering{arabic}

    % Optionale Einleitung und Struktur (Leitfaden-spezifisch)
    %\chapter*{Einleitung}
\addcontentsline{toc}{chapter}{Einleitung}
\label{chap:einleitung}

Die scientific community – also die Gemeinschaft der Menschen, die forschen und ihre Erkenntnisse mit anderen teilen wollen – steht allen offen. Es gibt keine formale, exkludierende Zugangsbarriere.

Spätestens seit Ihrer Immatrikulation an der Rheinischen Hochschule sind Sie Teil dieser scientifc community. Herzlich willkommen!

Ihr Studium bereitet Sie nicht erst aufs Forschen vor. Es steht Ihnen frei, vom ersten Tag an Ihre Zeit an der Rheinischen Hochschule forschend zu nutzen. Das gilt nicht nur für Seminare oder Projekte. Sie können auch schon Ihre Haus- und Abschlussarbeiten als Forschung konzipieren. Es spielt dafür keine Rolle, ob Sie Ihr Forschungsinteresse auf ein äußerst kleines Feld eingrenzen. Jeder Erkenntnisgewinn ist ein Dienst an der Wissenschaft, mag er Ihnen auch gering erscheinen.

Mit der Lektüre dieses Leitfaden ersetzen Sie selbstverständlich nicht Ihre aktive Teilnahme am Modul Wissenschaftliches Arbeiten. Es handelt sich hier lediglich um einige grundlegende Hinweise von übergeordnetem, die gesamte Hochschule betreffenden Charakter. Damit sollen insbesondere solche Aspekte des wissenschaftlichen Arbeitens geklärt werden, die sich verallgemeinern lassen.

Am Ende ist die Absprache mit Ihren Prüfern verbindlich, weil die individuelle Durchführung von Forschung und Lehre durch Artikel 5 Absatz 3 des Grundgesetzes geschützt ist.

Wissenschaftliches Arbeiten ist keine Tätigkeit, der Sie sich zuwenden, wenn Sie alles andere im Studium erledigt haben. Im Kern ist das Schreiben eines wissenschaftlichen Textes die Forschung selbst. Denn Forschung zielt auf Erkenntnisgewinn, den Sie nicht für sich behalten, sondern durch Veröffentlichung mit Ihrer Fachöffentlichkeit teilen.

Dafür müssen Sie Ihre Erkenntnisse in eine schriftliche Form bringen. Und um Ihr Publikum zu überzeugen, müssen Sie Ihre Aussagen intersubjektiv nachvollziehbar formulieren und sich auf vertrauenswürdige Quellen stützen. Beide Aspekte, die schriftliche Form und Ihre Nachweisführung, werden in diesem vorliegenden Leitfaden sowie in der separaten Praktische[n] Anleitung: Zitieren in wissenschaftlichen Arbeiten der Rheinischen Hochschule behandelt.

Was für Wissenschaft grundsätzlich gilt, das trifft auch auf diesen Leitfaden zu: Wir sind alle nur Menschen, die sich irren können und denen versehentliche Fehler unterlaufen. Falls Ihnen also etwas auffällt, das Ihnen fragwürdig oder offensichtlich falsch erscheint und deshalb geändert werden sollte, freuen wir uns jederzeit über Ihren Hinweis an unser Qualitätsmanagement: qm@rh-koeln.de. Vielen Dank für Ihre Hilfe!

\chapter*{Struktur und Gültigkeit dieses Leitfadens}
\addcontentsline{toc}{chapter}{Struktur und Gültigkeit dieses Leitfadens}
\label{chap:leitfaden-struktur}

Mit jeder Hausarbeit und Abschlussarbeit führen Sie vor, dass Sie das Wissen, welches Sie sich im Studium angeeignet haben, anwenden und praxisrelevante Fragen sachgerecht in einer wissenschaftlichen Methode beantworten können. Inwieweit Ihnen dies gelungen ist, unterliegt der Bewertung durch Ihre Prüfer.

Für die Bewertung wissenschaftlicher Arbeiten verwendet die Rheinische Hochschule ein einheitliches Schema. Es enthält sechs Abschnitte mit detaillierten Hinweisen und bestimmt, wieviele Punkte Sie für welche Leistungen erhalten.

\begin{table}[H]
\centering
\label{tab:leistungsabschnitte}
    \begin{tabular}{|l|l|c|}
    \specialrule{1.5pt}{0pt}{0pt}
    Nr. & Leistungsabschnitt & Punkte \\
    \specialrule{1.5pt}{0pt}{0pt}
    1. & Thema und Forschungsfrage bzw. Aufgabenstellung, Hypothesen & 10 \\
    \hline
    2. & Strukturierung & 10 \\
    \hline
    3. & Konzeptionell-theoretische Grundlagen & 20 \\
    \hline
    4. & Durchführung & 40 \\
    \hline
    5. & Quellen und Zitierweise & 10 \\
    \hline
    6. & Form und Stil & 10 \\
    \hline
    \end{tabular}
\caption[Übersicht der Leistungsabschnitte für die Bewertung von Haus- und Abschlussarbeiten an der Rheinischen Hochschule Köln.]{Übersicht der Leistungsabschnitte für die Bewertung von Haus- und Abschlussarbeiten an der Rheinischen Hochschule Köln.}
\end{table}

Damit Sie genau nachvollziehen können, welche Leistungen von Ihnen erwartet werden (Kategorien) und wie diese bewertet werden (Kriterien), folgt dieser Leitfaden dem Bewertungsschema.

In Anhang \ref{app:wirtschaftspsychologie-mediendesign} finden Sie ergänzende Hinweise für die Abschlussarbeiten in den Studiengängen Mediendesign und Wirtschaftspsychologie.

Der vorliegende Leitfaden ergänzt die Prüfungsordnungen (RSPO und SPO), die im Zweifelsfall verbindlich sind.
    %\newpage

    % Kapitel einbinden
    \chapter{Einleitung}
\label{chap:einleitung}

Die zunehmende Automatisierung verteilter Cloud- und Edge-Systeme stellt hohe Anforderungen an die Zuverlässigkeit der zugrundeliegenden Bereitstellungsmechanismen. Dieses einleitende Kapitel legt die Motivation der Untersuchung dar, ordnet das Vorhaben in den aktuellen Stand der Forschung ein, leitet daraus die zentrale Forschungsfrage sowie die Zielsetzung her und gibt einen Überblick über den Aufbau der Arbeit.

\section{Motivation}
\label{sec:motivation}

Die Ansprüche an heutige Cloud-Native-Infrastrukturen steigen stetig, vor allem an betriebsnahe Eigenschaften wie Ausfallsicherheit, automatisierte Selbstheilung und deterministische Reproduzierbarkeit. Genau für diese Metriken bilden Container-Orchestrierungssysteme wie Kubernetes in Verbindung mit deklarativen Bereitstellungsparadigmen das technische Fundament.\autocite{burns_borg_2016,burns_kubernetes_2022} Wurden Anwendungs- und Infrastrukturänderungen in klassischen \gl{ci}-/\gl{cd}-Pipelines historisch überwiegend über imperative, von außen angestoßene Deployment-Skripte nach dem Push-Prinzip in Zielumgebungen übertragen, hat sich mit GitOps ein gegenläufiges Betriebsmodell etabliert. Hierbei fungiert ein versionsverwaltetes Git-Repository als einzige verlässliche Wahrheitsquelle (\textit{\gl{ssot}}), während spezialisierte Software-Agenten innerhalb des Clusters den realen Systemzustand kontinuierlich mit dem deklarierten Soll-Zustand abgleichen.\autocite{beetz_gitops_2022,opengitops_principles_2021}

Aufgrund der zunehmenden Verbreitung leichtgewichtiger Kubernetes-Distributionen wie K3s verlagert sich der Einsatz solcher GitOps-Control-Planes zunehmend aus hochverfügbaren Rechenzentren in ressourcenbeschränkte Edge-, Zweigstellen- und Cloud-Instanzen.\autocite{k3s_architecture_2026,boehm_k3s_2021,koziolek_k3s_2023} Dies zeigt sich sowohl beim Betrieb verteilter \gl{iot}- und Edge-Knoten als auch bei kostenoptimierten Single-Node- und Kleinst-Clustern, die über öffentliche \gl{wan}-Strecken an externe Versionsverwaltungssysteme angebunden sind. In diesen Topologien kann eine unterbrechungsfreie Netzwerkverbindung zwischen dem Cluster-internen GitOps-Controller und dem externen Git-Repository nicht als dauerhaft gegeben vorausgesetzt werden. Außerdem gewinnt die Frage nach der Ressourcenstabilität der Control Plane auf Systemen mit strikt limitiertem Hauptspeicher (\gl{ram}) immer mehr an Bedeutung, da wiederholte Synchronisationsversuche und Timeouts bei Verbindungsstörungen zu Lastspitzen führen können. Diese Verschiebung der Einsatzgebiete rückt das dynamische Verhalten von GitOps-Operatoren unter transienter Netzwerkdegradation in den Fokus der Forschung und Systemarchitektur.\autocite{fogli_k3s_2021,ghosh_gitops_2025,basiri_chaos_2016}

\section{Stand der Forschung und Abgrenzung}
\label{sec:stand-der-forschung}

Die wissenschaftliche Literatur untersucht die Leistungsfähigkeit und Fehlertoleranz Cloud-nativer Orchestrierungs- und Bereitstellungssysteme bislang überwiegend in drei getrennten Teilbereichen. Einen ersten Schwerpunkt bildet die Evaluierung leichtgewichtiger Kubernetes-Distributionen auf ressourcenbeschränkten Knoten. Böhm und Wirtz profilieren den Ressourcenbedarf von MicroK8s und K3s im Vergleich zu einer regulären Kubernetes-Installation entlang typischer Lebenszyklusphasen und weisen für K3s einen deutlich geringeren Haupt- und Massenspeicherbedarf der Basisschicht nach.\autocite{boehm_k3s_2021} Koziolek und Eskandani erweitern diesen Vergleich auf MicroK8s, K3s, k0s und MicroShift und quantifizieren sowohl die minimale Speicherauslastung im Leerlauf als auch den Durchsatz der Control Plane unter Last.\autocite{koziolek_k3s_2023} Beide Studien setzen jedoch durchgängig störungsfreie lokale Netzwerke voraus und betrachten keine zusätzlichen GitOps-Operatoren. Fogli et\,al. untersuchen Standard-Kubernetes, K3s und KubeEdge zwar gezielt unter Paketverlusten, erhöhten Latenzen und Bandbreitenbegrenzungen in taktischen Edge-Netzwerken, richten den Fokus jedoch ausschließlich auf die Cluster-interne Verbindung zwischen Control Plane und entfernten Worker-Knoten.\autocite{fogli_k3s_2021}

Ein zweiter Forschungsstrang befasst sich mit dem konzeptionellen und empirischen Vergleich deklarativer GitOps-Werkzeuge. Aufbauend auf der architektonischen Einordnung des Pull-Prinzips durch Beetz und Harrer legen Ghosh et\,al. einen reproduzierbaren Leistungsbenchmark für die Operatoren ArgoCD, Flux\,CD und ConfigSync im Kontext von \textit{Intent-Based Networking} vor.\autocite{beetz_gitops_2022,ghosh_gitops_2025} Die Autoren erfassen die Verarbeitungs- und Bereitstellungslatenz, den Determinismus der Synchronisation sowie den Prozessor- und Speicherbedarf in Einzel- und Mehrfachszenarien.\autocite{ghosh_gitops_2025} Sämtliche Messreihen erfolgen allerdings unter idealen Netzwerkbedingungen, sodass das Konvergenzverhalten bei transienten Verbindungsstörungen zum externen Git-Repository sowie die resultierenden Wiederherstellungszeiten (\gl{mttr}) unberücksichtigt bleiben.

Drittens adressieren neuere Arbeiten der Systemforschung die automatisierte Zuverlässigkeitsprüfung von Kubernetes-Controllern und Operatoren. Sun et\,al. entwickeln mit \textit{Sieve} ein Testverfahren, das gezielt verzögerte Zustandsansichten, verworfene Ereignisse und asynchrone Teilausfälle zwischen Controller und \texttt{kube-apiserver} injiziert, um Konsistenz- und Konvergenzfehler in Reconciliation-Schleifen offenzulegen.\autocite{sun_sieve_2022} Daran anknüpfend überprüft \textit{Acto} von Gu et\,al. systematisch Zustandsübergänge von Kubernetes-Operatoren auf fehlerhafte Endzustände.\autocite{gu_acto_2023} Beide Ansätze zielen jedoch auf die qualitative Identifikation logischer Implementierungsfehler an der Cluster-internen \gl{api}-Grenze ab und erfassen weder reale Transportschicht-Störungen auf ausgehenden \gl{wan}-Strecken noch hochaufgelöste Latenz- und Ressourcenzeitreihen.

In der Zusammenschau bleibt die Schnittstelle zwischen ressourcenbeschränkten K3s-Umgebungen, deklarativen GitOps-Control-Planes und transienter Netzwerkdegradation empirisch weitgehend offen. Insbesondere fehlt eine quantitative Untersuchung, die den \textit{ArgoCD Application Controller} und den \textit{Repository-Server} selbst als \gl{sut} isoliert und die Kausalkette von probabilistischen Paketverlusten auf ausgehenden \gl{https}-Verbindungen (\gl{tcp}-Port~443) über Kernel-Retransmissionen und \gl{grpc}-Timeouts bis hin zur Reconcile-Latenz und Post-Fault-\gl{mttr} vermisst. Ebenso bleibt in bestehenden Benchmarks ungeklärt, ob wiederholte Synchronisationsversuche und blockierte Netzwerkroutinen auf speicherlimitierten Knoten zu Auslagerungsprozessen auf den Sekundärspeicher führen, welche Latenzmessungen durch Paging-Artefakte verfälschen würden.\autocite{gregg_systems_2020}

\section{Forschungsfrage und Zielsetzung}
\label{sec:forschungsfrage-zielsetzung}

Aus dieser Forschungslücke ergibt sich die Frage nach der quantifizierbaren Widerstandsfähigkeit (Resilienz) und dem Konvergenzverhalten moderner GitOps-Controller unter realen Fehlerbedingungen. Die vorliegende Arbeit widmet sich daher der folgenden Forschungsfrage:

\begin{quote}
\textit{Wie wirken sich transiente Paketverluste auf der Netzwerkverbindung zum externen Git-Repository auf die Reconcile-Latenz, die Zustandsübergänge, die Mean Time to Recovery (\gl{mttr}) sowie die Ressourcenauslastung des ArgoCD Application Controllers in einer ressourcenbeschränkten K3s-Umgebung aus?}
\end{quote}

Bei der fundierten Analyse dieser Forschungsfrage wird die Untersuchung nicht ausschließlich auf der Ebene einer qualitativen Funktionsprüfung geführt. Die Bewertung erfolgt gezielt aus der Perspektive der geschlossenen Regelungsschleife (\textit{Reconciliation Loop}), der zeitlichen Verteilung der Synchronisationslatenzen (Median und 95.-Perzentil), der Verweildauer in fehlerhaften Systemzuständen sowie der resultierenden \gl{cpu}-, Speicher- und Kernel-Paging-Auslastung.\autocite{beyer_sre_2016,nygard_release_2018} Hierdurch soll objektiv und reproduzierbar untersucht werden, inwiefern der \textit{ArgoCD Application Controller} (Version~3.6) bei massiven transienten \gl{tcp}-Paketverlusten seine Selbstheilungsfähigkeit wahrt, welche Verzögerungen durch exponentielle Backoff-Mechanismen und Kernel-Retransmissionen entstehen und ob die Stabilität des Gesamtsystems ohne Auslagerungsprozesse auf den Sekundärspeicher (\textit{Swap-Reinheit}) gewährleistet bleibt.\autocite{argocd_architecture_2026,gregg_systems_2020}

\section{Aufbau der Arbeit}
\label{sec:aufbau-der-arbeit}

Zur Beantwortung der Forschungsfrage ist die Arbeit in sechs Kapitel strukturiert. In Kapitel~\ref{chap:theoretische-grundlagen} werden die theoretischen Grundlagen deklarativer Infrastrukturmodelle, die Prinzipien des GitOps-Paradigmas, die regelungstheoretische Funktionsweise des Kubernetes-Operator-Patterns sowie die methodischen Grundlagen des Chaos Engineerings und der Resilienzmetriken erläutert.

Kapitel~\ref{chap:architektur-k3s-argocd} analysiert anschließend die spezifische Systemarchitektur der untersuchten Plattformkomponenten. Der Fokus liegt hierbei auf der leichtgewichtigen Kubernetes-Distribution K3s sowie der internen Mikroarchitektur, Zustandsmaschine und Fehlerbehandlung von ArgoCD im Kontext von Netzwerkstörungen und strikten Speichergrenzen.

Das vierte Kapitel beschreibt die Konzeption und Implementierung des im Rahmen dieser Arbeit entwickelten, vollständig als Code definierten Messframeworks (\textit{Infra-Chaos-Lab}). Dabei wird dargelegt, wie die ephemere Testumgebung auf Basis von Terraform und \gl{aws} \gl{ec2} automatisiert bereitgestellt, mit K3s, ArgoCD und Prometheus instrumentiert und um eine deterministische Fehlerinjektion auf Linux-Kernel-Ebene (\texttt{iptables}) sowie eine automatisierte Prüfung der Swap-Reinheit ergänzt wurde. Das fünfte Kapitel präsentiert und diskutiert die empirischen Messergebnisse der durchgeführten Chaos-Experimente anhand hochaufgelöster Zeitreihen- und Verteilungsanalysen. Die Arbeit endet im sechsten Kapitel mit einem konkreten Fazit, das die Forschungsfrage auf Basis der empirischen Befunde beantwortet, methodische Grenzen der Untersuchung reflektiert und einen Ausblick auf zukünftige Entwicklungen gibt.

    \chapter{Theoretische Grundlagen}
\label{chap:theoretische-grundlagen}

Die empirische Bewertung der Resilienz einer GitOps-Control-Plane erfordert zunächst eine theoretische Grundlage. Diese wird im Folgenden geschaffen, indem auf die Entwicklung von \gl{iac} hin zu deklarativen Zustandsmodellen, die formale Definition des GitOps-Betriebsmodells, die regelungstheoretische Funktionsweise der Reconciliation-Schleife sowie die Konzepte der Fehlertoleranz und des Chaos Engineerings eingegangen wird. Ein Blick auf die zugrundeliegenden Netzwerk- und Betriebssystemmechanismen verdeutlicht zudem die Kausalkette zwischen physischen Paketverlusten und dem Verhalten auf Anwendungsebene.

\section{Deklarative Infrastruktur und das GitOps-Paradigma}
\label{sec:deklarative-infrastruktur-gitops}

Um die Bereitstellung und den Betrieb komplexer Softwaresysteme reproduzierbar zu gestalten, hat sich das Konzept von \textit{Infrastructure as Code} (\gl{iac}) als grundlegendes Paradigma etabliert. Dabei werden sämtliche Infrastruktur- und Konfigurationsartefakte in maschinenlesbaren Definitionsdateien beschrieben, anstatt Systeme manuell oder interaktiv zu konfigurieren.\autocite{morris_iac_2020} Innerhalb von \gl{iac} sowie im Konfigurationsmanagement wird auf konzeptioneller Ebene strikt zwischen imperativen und deklarativen Ansätzen differenziert. Beide Ansätze verfolgen das Ziel einer automatisierten Systembereitstellung, wählen dafür aber unterschiedliche Steuerungsprinzipien.\autocite{burns_kubernetes_2022,morris_iac_2020}

\subsection{Imperative Steuerung versus deklarative Zustandsmodelle}
\label{subsec:imperativ-vs-deklarativ}

Der imperative Ansatz (prozedurales Modell) definiert eine explizite Abfolge von ausführbaren Einzelschritten, die ein System von seinem aktuellen Zustand in einen gewünschten Zielzustand überführen sollen. Klassische \gl{ci}-/\gl{cd}-Pipelines implementieren dieses Prinzip häufig über ein sogenanntes \textit{Push-Modell}: Nach einem erfolgreichen Build-Prozess führt ein externer Pipeline-Runner Kommandos oder Skripte (etwa \texttt{kubectl apply} oder \gl{ssh}-Befehle) direkt gegen die Zielumgebung aus.\autocite{beetz_gitops_2022,morris_iac_2020} Obwohl dieses Vorgehen bei einfachen Bereitstellungen intuitiv erscheint, weist es in verteilten Systemen strukturelle Schwächen auf. Da die Ausführung nur zum Zeitpunkt des Pipeline-Laufs stattfindet, bleiben nachträgliche manuelle Eingriffe oder schleichende Zustandsabweichungen im Zielsystem (\textit{Configuration Drift}) bis zum nächsten Deployment unbemerkt. Zudem erfordert das Push-Modell, dass weitreichende Schreib- und Administrationsberechtigungen für den Produktionscluster außerhalb der eigentlichen Laufzeitumgebung im \gl{ci}-System hinterlegt werden, was die Angriffsfläche der Infrastruktur vergrößert.\autocite{beetz_gitops_2022,opengitops_principles_2021}

Demgegenüber beschreibt das deklarative Modell ausschließlich den angestrebten Soll-Zustand (\textit{Desired State}) eines Systems in Form unveränderlicher Datenstrukturen (beispielsweise als \gl{yaml}-Manifeste), ohne die konkreten Ausführungsschritte zur Erreichung dieses Zustands vorzugeben.\autocite{burns_borg_2016,burns_kubernetes_2022} Die Verantwortung dafür, die Differenz zwischen dem aktuell beobachteten Ist-Zustand (\textit{Actual State} bzw. \textit{Live State}) und dem deklarierten Soll-Zustand zu berechnen und auszugleichen, wird vollständig an eine Laufzeitkomponente innerhalb des Zielsystems delegiert. Diese konsequente Trennung von Zustandsbeschreibung und Zustandsüberführung reduziert die Komplexität der Konfigurationsartefakte und bildet die Voraussetzung für selbstheilende Infrastrukturen.\autocite{burns_borg_2016,beyer_sre_2016}

\subsection{Formale Prinzipien des GitOps-Betriebsmodells}
\label{subsec:formale-prinzipien-gitops}

Das GitOps-Paradigma führt den deklarativen Infrastrukturansatz weiter und überträgt die Mechanismen der verteilten Versionsverwaltung auf den operativen Betrieb von Cloud-Native-Systemen. Um den Begriff wissenschaftlich und industriell eindeutig abzugrenzen, hat die \textit{OpenGitOps}-Arbeitsgruppe der \gl{cncf} vier verbindliche Kernprinzipien formuliert, die ein System erfüllen muss, um als GitOps-Architektur klassifiziert zu werden:\autocite{opengitops_principles_2021}

\begin{enumerate}
\item \textbf{Deklarativität (\textit{Declarative}):} Ein System, das über GitOps verwaltet wird, muss seinen gewünschten Zielzustand vollständig deklarativ ausdrücken.
\item \textbf{Versioniert und unveränderlich (\textit{Versioned and Immutable}):} Der gewünschte Zustand wird so gespeichert, dass Unveränderbarkeit, Versionierung und eine vollständige Versionshistorie erzwungen werden. In der Praxis übernimmt ein Git-Repository diese Funktion, wodurch jede Zustandsänderung durch einen kryptografisch signierbaren Commit-Hash (\gl{sha}) eindeutig identifizierbar und über ein \texttt{git revert} deterministisch umkehrbar ist.\autocite{beetz_gitops_2022,opengitops_principles_2021}
\item \textbf{Automatisch abgerufen (\textit{Pulled Automatically}):} Software-Agenten rufen die gewünschten Zustandserklärungen eigenständig und kontinuierlich aus der Datenquelle ab (\textit{Pull-Prinzip}). Im Gegensatz zum imperativen Push-Modell initiiert das Zielsystem die Verbindung nach außen, sodass keine Cluster-Zugangsdaten in externen Systemen exponiert werden müssen.\autocite{beetz_gitops_2022,opengitops_principles_2021}
\item \textbf{Kontinuierlich abgeglichen (\textit{Continuously Reconciled}):} Software-Agenten beobachten fortlaufend den tatsächlichen Systemzustand und unternehmen automatisierte Korrekturmaßnahmen, sobald eine Abweichung (\textit{Drift}) zum deklarierten Zustand festgestellt wird.\autocite{opengitops_principles_2021}
\end{enumerate}

Für die Untersuchung von Fehlerszenarien ist insbesondere das Zusammenspiel des dritten und vierten Prinzips von zentraler Bedeutung. Da der abgleichende Agent innerhalb des Clusters lokalisiert ist, die maßgebliche Wahrheitsquelle (das Git-Repository) jedoch außerhalb des Clusters liegt, entsteht eine permanente Abhängigkeit von der Netzwerkverbindung zwischen Control Plane und Versionsverwaltung.\autocite{beetz_gitops_2022,opengitops_principles_2021,ghosh_gitops_2025}

\section{Kontrolltheoretische Grundlagen: Das Operator-Pattern und die Reconciliation-Schleife}
\label{sec:operator-pattern-reconciliation}

Um das dynamische Verhalten von Kubernetes und darauf aufbauenden GitOps-Controllern wie ArgoCD unter Störeinflüssen präzise zu beschreiben, lässt sich deren Architektur auf Konzepte der Regelungstheorie (\textit{Control Theory}) zurückführen.\autocite{burns_borg_2016,lim_control_2009} In der Regelungstechnik überwacht ein Regler (\textit{Controller}) kontinuierlich die Regelgröße eines Systems, vergleicht diese mit einer vorgegebenen Führungsgröße (\textit{Setpoint}) und wirkt über eine Stellgröße korrigierend auf die Regelstrecke ein, um den Einfluss externer Störgrößen (\textit{Disturbances}) zu kompensieren.\autocite{beyer_sre_2016,lim_control_2009}

Kubernetes überträgt dieses Prinzip in Form der \textit{Reconciliation-Schleife} (\textit{Reconcile Loop}) auf die Softwaretechnik. Jeder Controller innerhalb der Kubernetes Control Plane arbeitet als geschlossener Regelkreis, der nach dem Muster \textit{Observe -- Diff -- Act} (Beobachten -- Vergleichen -- Handeln) operiert:\autocite{burns_borg_2016,dobies_operators_2020}

\begin{equation}
\Delta S(t) = S_{\text{desired}}(t) - S_{\text{actual}}(t)
\label{eq:reconcile-diff}
\end{equation}

Solange die Zustandsdifferenz $\Delta S(t) = \emptyset$ beträgt, befindet sich das System im eingeschwungenen Gleichgewichtszustand (\textit{Steady State} bzw. \texttt{Synced} und \texttt{Healthy}). Tritt eine Abweichung $\Delta S(t) \neq \emptyset$ auf -- sei es durch eine Aktualisierung des Soll-Zustands $S_{\text{desired}}(t)$ im Git-Repository oder durch einen Ausfall bzw. eine manuelle Veränderung im Cluster $S_{\text{actual}}(t)$ --, berechnet der Controller die erforderlichen \gl{api}-Operationen, um $\Delta S(t)$ wieder gegen Null zu führen.\autocite{burns_borg_2016,dobies_operators_2020}

Ein wesentliches Architekturmerkmal dieser Regelungsschleife ist die \textit{pegelgetriebene Triggerung} (\textit{Level-Triggered Reconciliation}) im Kontrast zu einer rein \textit{flankengetriebenen Verarbeitung} (\textit{Edge-Triggered}).\autocite{burns_borg_2016,dobies_operators_2020} Während ein flankengetriebenes System lediglich auf den unmittelbaren Zeitpunkt eines Ereignisses (etwa einen Webhook bei einem \texttt{git push}) reagiert, bewertet ein pegelgetriebenes System periodisch den gesamten aktuellen Zustand im Vergleich zum Zielzustand. Geht ein einzelnes Benachrichtigungsereignis aufgrund einer transienten Netzwerkstörung verloren, stellt die pegelgetriebene Architektur sicher, dass das System bei der nächsten erfolgreichen Zustandsabfrage dennoch zuverlässig konvergiert.\autocite{burns_borg_2016,dobies_operators_2020}

Das \textit{Kubernetes-Operator-Pattern} erweitert diesen nativen Regelungsmechanismus, indem domänenspezifisches Betriebswissen in eigene Controller gegossen und über \textit{Custom Resource Definitions} (\gl{crd}) direkt in die Kubernetes-\gl{api} integriert wird.\autocite{dobies_operators_2020} ArgoCD nutzt dieses Muster, indem es die Custom Resource \texttt{Application} definiert. Diese Ressource verknüpft eine externe Quelle (Repository-URL, Ziel-Revision und Verzeichnispfad) mit einem Ziel-Cluster und spezifischen Synchronisationsrichtlinien (wie \texttt{selfHeal} zur automatischen Korrektur von Cluster-Drift und \texttt{prune} zum Entfernen veralteter Ressourcen).\autocite{argocd_architecture_2026,dobies_operators_2020}

\section{Fehlertoleranz, Resilienz und Chaos Engineering in verteilten Systemen}
\label{sec:fehlertoleranz-resilienz-chaos}

Verteilte Systeme, die über Netzwerkgrenzen hinweg kommunizieren, unterliegen zwangsläufig partiellen Ausfällen. Um die Zuverlässigkeit solcher Architekturen wissenschaftlich zu bewerten, bedarf es einer klaren Abgrenzung der Begriffe aus der Verlässlichkeitsforschung (\textit{Dependability}) sowie etablierter Metriken zur Quantifizierung des Wiederherstellungsverhaltens.

\subsection{Fehlerkette (Fault -- Error -- Failure) und Resilienzbegriff}
\label{subsec:fehlerkette-resilienz}

Nach der grundlegenden Taxonomie von Avizienis et\,al. wird die Entstehung von Systemausfällen als dreistufige Kausalkette modelliert:\autocite{avizienis_dependable_2004} Ein \textit{Fault} (Fehlerursache bzw. Störung) bezeichnet die zugrundeliegende oder externe Ursache eines Fehlverhaltens, im Kontext dieser Arbeit etwa den partiellen Verlust von Netzwerkpaketen auf der Transportschicht. Führt dieser \textit{Fault} innerhalb des Systems zu einem abweichenden internen Teilzustand -- beispielsweise einem abgebrochenen Git-Abruf oder dem Übergang der ArgoCD-Zustandsmaschine in die Bedingung \texttt{ComparisonError} --, liegt ein \textit{Error} (Fehlzustand) vor. Erst wenn dieser interne Fehlzustand die nach außen sichtbare Dienstleistung des Gesamtsystems beeinträchtigt oder vollständig unterbricht, spricht man von einem \textit{Failure} (Ausfall).\autocite{avizienis_dependable_2004}

Der Begriff der \textit{Resilienz} (\textit{Resilience}) beschreibt vor diesem Hintergrund die Fähigkeit eines Systems, trotz des Auftretens von \textit{Faults} und resultierenden \textit{Errors} ein akzeptables Dienstniveau aufrechtzuerhalten (\textit{Graceful Degradation}) und nach dem Abklingen der Störung eigenständig sowie innerhalb einer definierten Zeitspanne wieder in den regulären Betriebszustand zurückzukehren.\autocite{beyer_sre_2016,avizienis_dependable_2004} In einem GitOps-System manifestiert sich Resilienz auf zwei Ebenen: Erstens darf ein Kommunikationsverlust zum externen Git-Repository die bereits im Cluster laufenden Produktiv-Workloads (\textit{Data Plane}) nicht beeinträchtigen (\textit{statische Resilienz}). Zweitens muss die \textit{Control Plane} nach Ende der Netzwerkstörung ohne manuellen Eingriff und ohne Ressourcenerschöpfung (wie \gl{oom}-Abstürze) wieder in den Zustand \texttt{Synced} und \texttt{Healthy} konvergieren (\textit{dynamische Resilienz}).\autocite{fogli_k3s_2021,sun_sieve_2022,argocd_architecture_2026}

\subsection{Zeitliche Resilienzmetriken: Verweildauer und Mean Time to Recovery (MTTR)}
\label{subsec:zeitliche-resilienzmetriken}

Zur quantitativen Erfassung der dynamischen Resilienz werden in der Systemevaluierung und im \textit{Site Reliability Engineering} (\gl{sre}) zeitbasierte Kennzahlen herangezogen.\autocite{beyer_sre_2016,nygard_release_2018} Im Rahmen einer kontrollierten Fehlerinjektion, die zum Zeitpunkt $t_{\text{fault\_start}}$ beginnt und zum Zeitpunkt $t_{\text{fault\_end}}$ endet, lassen sich zwei präzise abgegrenzte Zeitintervalle definieren:

\begin{enumerate}
\item \textbf{Transiente Verweildauer (\textit{Dwell Time} $T_{\text{dwell}}$):} Bezeichnet die gesamte Zeitspanne zwischen dem erstmaligen Eintritt des Systems in einen degradierten oder fehlerbehafteten Zustand ($t_{\text{degraded}}$) und der vollständigen Wiederherstellung des regulären Systemzustands ($t_{\text{recovered}}$):
\begin{equation}
T_{\text{dwell}} = t_{\text{recovered}} - t_{\text{degraded}}
\label{eq:dwell-time}
\end{equation}
\item \textbf{Mean Time to Recovery ($\text{\gl{mttr}}_{\text{post-fault}}$):} Im engeren messtechnischen Sinn eines kontrollierten Störungsexperiments beschreibt die Erholungszeit nach Störungsende die Differenz zwischen dem Zeitpunkt, an dem die externe Netzwerkstörung behoben ist ($t_{\text{fault\_end}}$), und dem Zeitpunkt, an dem die Control Plane den fehlerfreien Zielzustand (\texttt{Synced} und \texttt{Healthy} ohne Fehler-Conditions) wieder vollständig erreicht hat:\autocite{beyer_sre_2016,nygard_release_2018}
\begin{equation}
\text{\gl{mttr}}_{\text{post-fault}} = t_{\text{recovered}} - t_{\text{fault\_end}}
\label{eq:mttr}
\end{equation}
\end{enumerate}

Diese Differenzierung ist für die Bewertung von GitOps-Controllern essenziell: Während $T_{\text{dwell}}$ auch durch die Dauer der von außen aufgeprägten Netzwerkstörung bestimmt wird, isoliert $\text{\gl{mttr}}_{\text{post-fault}}$ die intrinsische Erholungsgeschwindigkeit der Control Plane, welche maßgeblich durch Polling-Intervalle, Timeout-Schwellenwerte und exponentielle Backoff-Algorithmen determiniert wird.\autocite{beyer_sre_2016,argocd_architecture_2026}

\subsection{Methodik des Chaos Engineerings}
\label{subsec:methodik-chaos-engineering}

Um die genannten Resilienzmetriken nicht erst bei unvorhersehbaren Produktionsausfällen passiv zu beobachten, sondern systematisch im Experiment zu bestimmen, wurde die Disziplin des \textit{Chaos Engineerings} entwickelt. Nach Basiri et\,al. bezeichnet Chaos Engineering das Durchführen gezielter, kontrollierter Experimente an einem verteilten System, um Vertrauen in dessen Fähigkeit aufzubauen, turbulenten und fehlerhaften Bedingungen im Betrieb standzuhalten.\autocite{basiri_chaos_2016}

Ein wissenschaftlich fundiertes Chaos-Experiment folgt dabei vier streng strukturierten Phasen:\autocite{basiri_chaos_2016} Zunächst wird der \textit{Steady State} (Normalzustand) des Systems anhand messbarer Ausgabemetriken (hier: Reconcile-Latenz, Synchronisationsstatus und Ressourcenverbrauch) als Referenzbasis (\textit{Baseline}) definiert. Darauf aufbauend wird die Hypothese formuliert, dass dieser Steady State auch bei Eintritt einer spezifischen Störung erhalten bleibt beziehungsweise nach deren Ende innerhalb definierter Grenzen wiederhergestellt wird. Anschließend werden reale Fehlervariablen -- wie Netzwerk-Paketverluste, Latenzen oder Serverabstürze -- kontrolliert in das System injiziert, während der \textit{Blast Radius} (Wirkungsbereich der Störung) strikt begrenzt und überwacht wird. Abschließend wird die Abweichung zwischen dem Verhalten unter Störung und der Baseline statistisch ausgewertet, um Schwachstellen in der Konvergenzlogik oder der Ressourcendimensionierung aufzudecken.\autocite{basiri_chaos_2016}

\section{Netzwerkschicht-Interaktionen und Betriebssystem-Einflussfaktoren}
\label{sec:netzwerkschicht-betriebssystem}

Da die Kommunikation zwischen einem GitOps-Controller und einem externen Git-Repository (z.\,B. GitHub) über \gl{https} (\gl{tcp}-Port 443) abgewickelt wird, wirkt sich ein Paketverlust auf der Netzwerkschicht unmittelbar auf das Zeitverhalten der höheren Protokollschichten aus. Gemäß der \gl{tcp}-Spezifikation (\gl{rfc} 9293) sichert das \textit{Transmission Control Protocol} eine zuverlässige Datenübertragung durch Sequenznummern, Empfangsbestätigungen (\gl{ack}) und zeitgesteuerte Neuübertragungen (\textit{Retransmissions}) ab.\autocite{eddy_rfc9293_2022}

Gehen Netzwerkpakete -- etwa durch Überlastung auf \gl{wan}-Strecken oder durch eine gezielte Injektion mittels des Linux-Kernel-Subsystems \texttt{netfilter}/\texttt{iptables} -- mit einer hohen Wahrscheinlichkeit $p$ verloren, greifen im \gl{tcp}-Stack des Betriebssystems gestaffelte Erholungsmechanismen:\autocite{eddy_rfc9293_2022,netfilter_iptables_2026} Beim \textit{Three-Way-Handshake} zum Verbindungsaufbau (\texttt{SYN} $\rightarrow$ \texttt{SYN-ACK} $\rightarrow$ \texttt{ACK}) sowie bei vollständigem Ausbleiben von Bestätigungspaketen tritt ein \textit{Retransmission Timeout} (\gl{rto}) ein. Nach jedem erfolglosen Übertragungsversuch verdoppelt der Linux-Kernel das \gl{rto}-Intervall gemäß einem exponentiellen Backoff-Verfahren.\autocite{eddy_rfc9293_2022} Bei einer hohen Paketverlustrate von beispielsweise $p = 0{,}80$ (80\,\% Paketverlust) führt bereits die sequenzielle Verkettung aus \gl{dns}-Auflösung, \gl{tcp}-Handshake, \gl{tls}-1.3-Schlüsselaushandlung und dem eigentlichen Git-Protokolltransfer (\texttt{git ls-remote} bzw. \texttt{git fetch}) dazu, dass einzelne Abfragen über viele Sekunden blockieren oder das applikationsseitige Timeout überschreiten.\autocite{argocd_architecture_2026,eddy_rfc9293_2022}

Auf Seiten der Control Plane hat dieses Verhalten direkte Auswirkungen auf den Ressourcenbedarf. Während blockierte Netzwerk-Goroutinen im Controller und im Repository-Server auf den Ablauf von \gl{tcp}- und \gl{http}-Timeouts warten, belegen die zugehörigen Datenstrukturen, Puffer und \gl{tls}-Kontexte weiterhin Arbeitsspeicher (\textit{Working Set Memory}). Auf kostenoptimierten Cloud- oder Edge-Knoten mit knapp bemessenem physischem \gl{ram} birgt dies ein methodisches Risiko für die Leistungsbewertung: Überschreitet der Speicherbedarf die Grenze des physischen Hauptspeichers, beginnt der Linux-Kernel, Speicherseiten auf den Sekundärspeicher auszulagern (\textit{Swapping} bzw. \textit{Paging}), oder beendet Prozesse über den \gl{oom}-Killer.\autocite{gregg_systems_2020,kernel_vmstat_2026} Da Zugriffe auf ein \gl{ssd}-basiertes Swapfile um Größenordnungen höhere Latenzen aufweisen als Zugriffe auf physischen \gl{dram}, würde ein unbemerktes Swapping die gemessenen Reconcile-Latenzen verfälschen (\textit{Messartefakt}). Eine wissenschaftlich valide Latenzmessung in ressourcenbeschränkten Umgebungen setzt daher zwingend den Nachweis voraus, dass während des gesamten Experiments keinerlei Speicherseiten ein- oder ausgelagert wurden ($\Delta \text{pswpin} = 0$ und $\Delta \text{pswpout} = 0$ in \texttt{/proc/vmstat}).\autocite{gregg_systems_2020,kernel_vmstat_2026}

    \chapter{Architektur und Konvergenzmechanismen von K3s und ArgoCD}
\label{chap:architektur-k3s-argocd}

Das untersuchte Zielsystem besteht aus der leichtgewichtigen Kubernetes-Distribution K3s als Laufzeitbasis und ArgoCD als deklarativer GitOps-Control-Plane. Um Latenzspitzen, Zustandsübergänge und Speicherprofile im Experiment kausal einordnen zu können, werden im Folgenden die Komponentenstruktur sowie die internen Polling-, Timeout- und Begrenzungsmechanismen beider Plattformen analysiert.

\section{Leichtgewichtige Container-Orchestrierung mit K3s}
\label{sec:leichtgewichtige-orchestrierung-k3s}

Standard-Kubernetes verteilt Control-Plane-Komponenten wie \texttt{kube-apiserver}, \texttt{etcd}, \texttt{kube-scheduler} und \texttt{kube-controller-manager} auf getrennte Prozesse und Container.\autocite{burns_borg_2016} Auf kompakten Edge-Knoten oder kleinen Cloud-Instanzen (etwa \gl{aws} \texttt{c7i-flex.large} mit 2~\gl{vcpu}s und 4~GiB \gl{ram}) belegt diese Basisschicht bereits einen spürbaren Teil des verfügbaren Hauptspeichers, noch bevor GitOps-Operatoren oder Monitoring-Dienste gestartet werden.\autocite{koziolek_k3s_2023}

Die \gl{cncf}-zertifizierte Distribution \textit{K3s} konsolidiert diesen Aufbau, indem sie sämtliche Control-Plane- und Worker-Komponenten samt \texttt{containerd} und dem \gl{cni}-Plugin \texttt{Flannel} in einer einzigen Binärdatei bündelt und veraltete Cloud-Provider-Treiber entfernt.\autocite{k3s_architecture_2026} Für Single-Node-Cluster ersetzt K3s zudem den speicherintensiven Key-Value-Store \texttt{etcd} standardmäßig durch eine eingebettete SQLite-Datenbank über die Abstraktionsschicht \textit{Kine}, die \texttt{etcd}-Aufrufe transparent in relationale Transaktionen übersetzt.\autocite{k3s_architecture_2026}

Empirische Vergleiche zeigen dabei, dass die Vorteile nicht pauschal in jeder Dimension gleich groß sind: Während Böhm und Wirtz vor allem deutliche Einsparungen beim Massenspeicher und schnellere Startzeiten feststellen, aber nur moderate \gl{ram}-Unterschiede zum regulären Kubernetes messen,\autocite{boehm_k3s_2021} dokumentieren Koziolek und Eskandani im Vergleich mehrerer Lightweight-Distributionen eine kompakte Leerlauf-Speicherauslastung von K3s.\autocite{koziolek_k3s_2023} Für kontrollierte Experimente lässt sich der Eigenbedarf weiter senken, indem vorinstallierte Zusatzdienste wie \textit{Traefik} und der \textit{Metrics-Server} bei der Installation deaktiviert werden (\texttt{--disable traefik --disable metrics-server}). Die volle \gl{api}- und \gl{crd}-Kompatibilität bleibt erhalten, während nahezu das gesamte Knotenbudget dem \gl{sut} und der Metrikerfassung zur Verfügung steht.

\section{Komponentenarchitektur und Zustandsmaschine von ArgoCD}
\label{sec:komponentenarchitektur-argocd}

ArgoCD trennt seine Control Plane in drei zentrale Microservices und einen internen Cache.\autocite{argocd_architecture_2026} Der \textit{\gl{api}-/Web-Server} (\texttt{argocd-\allowbreak{}server}) bietet \gl{grpc}-, \gl{rest}- und Web-\gl{ui}-Schnittstellen für Administration, Authentifizierung und manuelle Sync-Befehle, übernimmt aber keine periodischen Hintergrundabgleiche. Für die externe Git-Kommunikation ist der zustandsfreie \textit{Repository-Server} (\texttt{argocd-\allowbreak{}repo-\allowbreak{}server}) zuständig: Er klont konfigurierte Repositories, ruft Änderungen ab und rendert deklarative Vorlagen wie Kustomize-Overlays oder Helm-Charts in konkrete Kubernetes-Manifeste.\autocite{argocd_architecture_2026,kubernetes_kustomize_2026} Das eigentliche Regelungsherz bildet der Go-basierte \textit{Application Controller} (\texttt{argocd-application-controller}). Er beobachtet alle \texttt{Application}-Ressourcen, bezieht über den \texttt{repo-server} den Soll-Zustand, vergleicht diesen mithilfe des internen \texttt{redis}-Caches mit dem Live-Zustand im Cluster und führt bei Abweichungen Korrekturoperationen gegen den \texttt{kube-apiserver} aus.\autocite{argocd_architecture_2026}

Den Zustand jeder Anwendung modelliert der Controller über zwei orthogonale Dimensionen sowie eine Liste von \textit{Conditions} (Tabelle~\ref{tab:argocd-zustandsdimensionen}).

\begin{table}[H]
    \centering
    \small
    \begin{tabularx}{\textwidth}{>{\raggedright\arraybackslash}p{3.9cm}|>{\raggedright\arraybackslash}p{3.5cm}|>{\raggedright\arraybackslash}X}
    \textbf{Zustandsdimension} & \textbf{Mögliche Werte} & \textbf{Bedeutung im Reconcile-Zyklus} \\
    \specialrule{1.5pt}{0pt}{0pt}
    Sync-Status\newline (\texttt{\footnotesize .status.sync.status}) & \texttt{Synced}, \texttt{OutOfSync}, \texttt{Unknown} & Zeigt an, ob der Live-Zustand im Cluster mit den Git-Manifesten übereinstimmt. Scheitert der Git-Abruf, wechselt der Status ggf.\ in \texttt{Unknown}. \\
    \hline
    Health-Status\newline (\texttt{\footnotesize .status.health.status}) & \texttt{Healthy}, \texttt{Progressing}, \texttt{Degraded}, \texttt{Suspended}, \texttt{Missing}, \texttt{Unknown} & Bewertet die betriebliche Gesundheit der untergeordneten Kubernetes-Objekte (z.\,B.\ laufende Pod-Replicas eines \texttt{Deployment}). \\
    \hline
    Conditions\newline (\texttt{\footnotesize .status.conditions}) & u.\,a.\ \texttt{ComparisonError}, \texttt{SyncError}, \texttt{InvalidSpecError} & Protokolliert Fehler während des Abgleichs. Bei Netzwerkverlust zum Git-Server hinterlegt der Controller eine \texttt{ComparisonError}-Bedingung. \\
    \hline\hline
    \end{tabularx}
\caption{Zustandsdimensionen einer ArgoCD-\texttt{Application}-Ressource, eigene Darstellung nach Argo Project (2026).}
\label{tab:argocd-zustandsdimensionen}
\end{table}

Diese Trennung begründet ein zentrales Resilienzmerkmal: Selbst wenn eine Netzwerkstörung zum Git-Server eine \texttt{ComparisonError}-Bedingung auslöst, bleiben die bereits laufenden Produktiv-Workloads im Cluster \texttt{Healthy}. Solange der Controller stabil weiterarbeitet, bleibt der Fehler auf die Synchronisationsfähigkeit der Control Plane begrenzt.

\section{Reconcile-Zyklus, Timeout-Kaskaden und Ressourcen-Constraints}
\label{sec:reconcile-zyklus-timeouts-ressourcen}

Im Normalbetrieb pollt ArgoCD Git-Repositories periodisch, um Netzwerk- und \gl{cpu}-Last gering zu halten. Der Standardwert von \texttt{timeout.\allowbreak{}reconciliation} liegt bei \texttt{180s}; das eingesetzte Helm-Chart~\texttt{10.18.2} setzt \texttt{120s} zuzüglich bis zu \texttt{60s} Jitter (\texttt{timeout.\allowbreak{}reconciliation.\allowbreak{}jitter}).\autocite{argocd_reconcile_2026,argo_helm_2026} Zusätzlich cacht der \texttt{argocd-\allowbreak{}repo-\allowbreak{}server} gerenderte Manifeste in Redis: Bleibt die Commit-\gl{sha} bei einer \texttt{git ls-remote}-Abfrage unverändert, entfällt ein vollständiger \texttt{git fetch}- und Kustomize-Durchlauf.\autocite{argocd_reconcile_2026} Für ein zweiminütiges Störungsexperiment folgt daraus eine zentrale Designanforderung: Ohne aktiven Eingriff könnte das Chaos-Fenster vollständig zwischen zwei Polling-Zyklen liegen oder aus dem Cache bedient werden. Erst die Annotation \texttt{argocd.\allowbreak{}argoproj.io/\allowbreak{}refresh=hard} invalidiert den Manifest-Cache und erzwingt eine unmittelbare Git-Abfrage über den \texttt{repo-server}.\autocite{argocd_reconcile_2026}

Trifft eine solche Abfrage auf eine paketverlustbehaftete Strecke, greifen zwei technisch getrennte Mechanismen:

In der \textit{Reconciliation-Phase} ruft der \texttt{application-\allowbreak{}controller} per \gl{grpc} den \texttt{repo-server} auf, der seinerseits einen Git-Prozess über \gl{https} (Port~443) startet. Verzögern \gl{tcp}-Retransmissionen diesen Abruf, steigen sowohl die Dauer der Git-Anfrage (\texttt{argocd\_\allowbreak{}git\_\allowbreak{}request\_\allowbreak{}duration\_\allowbreak{}seconds}) als auch die gesamte Reconcile-Latenz (\texttt{argocd\_\allowbreak{}app\_\allowbreak{}reconcile}).\autocite{argocd_metrics_2026} Bricht der Git-Abruf per Timeout ab, setzt der Controller eine \texttt{ComparisonError}-Condition und überführt den \texttt{SyncStatus} in \texttt{Unknown}, ohne überhaupt eine Sync-Operation (\texttt{argocd\_\allowbreak{}app\_\allowbreak{}sync\_\allowbreak{}total}) gegen Kubernetes zu starten. Weitere Abgleichversuche werden in diesem Zustand über \gl{tcp}-/\gl{grpc}-Timeouts und die interne Workqueue- bzw.\ Polling-Planung gesteuert.

Die applikationsseitige \textit{Retry-Richtlinie} (\texttt{syncPolicy.retry.backoff}) greift dagegen erst dann, wenn bei bekanntem Soll-Zustand und erkanntem Drift (\texttt{OutOfSync}) eine aktive Auto-Sync-Operation (\texttt{selfHeal: true}) gegen den \texttt{kube-apiserver} selbst fehlschlägt.\autocite{argocd_autosync_2026} Bei \texttt{limit: 5}, Basiswartezeit $d_0 = 5\,\text{s}$, Faktor $f = 2$ und Deckel $d_{\max} = 180\,\text{s}$ berechnet sich die Wartezeit $d_n$ vor dem $n$-ten Wiederholungsversuch ($n \in \{1, \dots, 5\}$) zu:

\begin{equation}
d_n = \min(d_0 \cdot f^{n-1},\; d_{\max})
\label{eq:exponential-backoff}
\end{equation}

Die resultierende Staffelung ($5\,\text{s}, 10\,\text{s}, 20\,\text{s}, 40\,\text{s}, 80\,\text{s}$) verhindert bei anhaltenden Sync-Fehlern einen \textit{Retry Storm}.\autocite[Kap.~22, S.~259--272]{beyer_sre_2016} Bei reinen Git-Netzwerkstörungen bestimmen jedoch primär ablaufende \gl{tcp}-/\gl{grpc}-Timeouts sowie der Abstand bis zum nächsten Reconcile-Zyklus die Erholungszeit $\text{\gl{mttr}}_{\text{post-fault}}$.

Gleichzeitig halten wartende Reconcile-Worker im Controller sowie blockierte \gl{tls}-Kontexte und Kustomize-Kindprozesse im Repo-Server Arbeitsspeicher gebunden. Auf ressourcenbeschränkten Knoten werden beide Pods daher mit engen Speicherlimits versehen (\texttt{512Mi} für den Controller, \texttt{384Mi} für den Repo-Server). Überschreitet ein Container sein Limit, beendet das Linux-Cgroups-Subsystem den Prozess sofort mit \texttt{OOMKilled} (Exit-Code~137).\autocite[Kap.~5]{burns_kubernetes_2022} Resilienz bemisst sich somit nicht nur daran, ob ArgoCD nach einer Störung überhaupt konvergiert, sondern am Zusammenspiel aus Latenzanstieg, $\text{\gl{mttr}}_{\text{post-fault}}$ und dem Sicherheitsabstand zur \gl{oom}-Grenze bei voller Swap-Reinheit.

    \chapter{Konzeption und Implementierung des automatisierten Messframeworks}
\label{chap:konzeption-implementierung}

Dieses Kapitel beschreibt Architektur und Implementierung des entwickelten Messframeworks (\textit{Infra-Chaos-Lab}). Erläutert werden die deklarative Bereitstellung der ephemeren Testumgebung auf \gl{aws} \gl{ec2}, die Kernel-seitige Fehlerinjektion via \texttt{iptables}, die hochauflösende Instrumentierung ohne blinde Flecken sowie die autarke On-Node-Versuchsausführung inklusive Kontrolle der Swap-Reinheit.

\section{Architektur und deklarative Bereitstellung der ephemeren Testumgebung}
\label{sec:architektur-bereitstellung-testumgebung}

Um Messreihen frei von Konfigurationsrückständen, fragmentierten Datenbankzuständen oder verwaisten Netzwerkregeln durchzuführen, folgt das Messframework dem Prinzip ephemerer Testumgebungen.\autocite{morris_iac_2020,brikman_terraform_2022} Die Infrastruktur wird für jede Testreihe per \gl{iac} neu erzeugt, mit fest gepinnten Softwareversionen provisioniert und nach Datenerhebung vollständig abgebaut.

\subsection{Cloud-Infrastruktur, Netzwerktopologie und Cloud-Init-Injektion}
\label{subsec:cloud-infrastruktur-cloud-init}

Terraform stellt in der \gl{aws}-Region \texttt{eu-central-1} eine isolierte \gl{vpc} (\texttt{10.0.0.0/16}) mit öffentlichem Subnetz (\texttt{10.0.1.0/24}) und Internet-Gateway bereit (\texttt{network.tf}).\autocite{brikman_terraform_2022,aws_ec2_2026} Zur Begrenzung externer Störeinflüsse ermittelt \texttt{security\_groups.tf} beim \texttt{terraform apply} über \texttt{checkip.amazonaws.com} die öffentliche \gl{wan}-IP des Versuchsleiters und beschränkt alle eingehenden Regeln auf das zugehörige \texttt{/24}-Subnetz (Ports \texttt{22}, \texttt{6443}, \texttt{80}, \texttt{443} und NodePort-Bereich \texttt{30000}--\texttt{32767}). Ausgehender Verkehr (\textit{Egress}) bleibt offen, damit Paketquellen, Container-Images und GitHub über Port \texttt{443} erreichbar sind.\autocite{aws_ec2_2026}

Als Rechenknoten (\texttt{compute.tf}) dient eine \gl{ec2}-Instanz vom Typ \texttt{t3.medium} (2~\gl{vcpu}s, 4~GiB physischer \gl{ram}) mit verschlüsseltem 30-GB-\texttt{gp3}-Root-Volume und Ubuntu 24.04 LTS (\textit{Noble Numbat}).\autocite{aws_ec2_2026,cloudinit_docs_2026} Da \gl{aws} das User-Data-Feld auf 16~KiB begrenzt, komprimiert \texttt{templatefile()} die Helm-Values für ArgoCD und Prometheus, das Baseline-Applikationsmanifest sowie das Bootstrap-Skript per \texttt{base64gzip()}. Cloud-Init entpackt diese Dateien unter \texttt{/opt/} bzw.\ \texttt{/usr/local/bin/} und startet den Bootstrap automatisch via \texttt{runcmd}. Über \texttt{user\_data\_replace\_on\_change = true} erzwingt jede Konfigurationsänderung eine frische Instanz.\autocite{brikman_terraform_2022,cloudinit_docs_2026}

\subsection{Automatisierter Cluster-Bootstrap und Ressourcenallokation}
\label{subsec:cluster-bootstrap-ressourcenallokation}

Das Skript \texttt{bootstrap-cluster.sh} richtet das System deterministisch mit explizit gepinnten Versionen ein: Zunächst erzeugt es ein 4-GiB-Swapfile (\texttt{/swapfile}), um \gl{oom}-Abbrüche während der parallelen Image-Entpackung im Bootstrap zu verhindern, setzt jedoch \texttt{vm.swappiness=1} (\texttt{/etc/sysctl.d/99-swappiness.conf}), damit der Kernel im späteren Experimentbetrieb keinen Speicher proaktiv auslagert.\autocite{gregg_systems_2020,kernel_vmstat_2026} Danach bezieht das Skript die öffentliche IPv4-Adresse über \gl{imds} (\texttt{--tls-san}) und installiert K3s in der gepinnten Version \texttt{v1.37.2+k3s1} ohne Traefik und Metrics-Server (\texttt{--disable traefik --disable metrics-server --kubelet-arg fail-swap-on=false}).\autocite{k3s_architecture_2026,aws_ec2_2026}

Nach der Installation von Helm~v3 werden die Namespaces \texttt{argocd}, \texttt{monitoring} und \texttt{workloads} priorisiert bestückt:\autocite{burns_kubernetes_2022,helm_docs_2026}

\begin{enumerate}
\item \textbf{ArgoCD Control Plane (Priorität 1):} Über das Chart \texttt{argo/argo-cd} (Chart-Version \texttt{10.18.2}, ArgoCD~\texttt{v3.6.2}) werden ausschließlich die Kernkomponenten ausgerollt; \texttt{dex} und \texttt{notifications} sind deaktiviert sowie \texttt{applicationSet} auf null Replicas skaliert.\autocite{argocd_architecture_2026} Application Controller (Port \texttt{8082}) und Repo-Server (Port \texttt{8084}) exponieren ihre Prometheus-Endpunkte über Pod-Annotationen.
\item \textbf{Schlanker Monitoring-Stack (Priorität 2):} Über \texttt{prometheus-community/\allowbreak{}prometheus} (Chart-Version \texttt{29.48.0}, Prometheus~\texttt{v3.16.0}) läuft ein einzelner Prometheus-Server samt \texttt{prometheus-node-exporter} und \texttt{kube-state-metrics}; auf Alertmanager, Pushgateway und ein separates Grafana-Deployment wird zur Minimierung des Observer-Overheads verzichtet.\autocite{brazil_prometheus_2018,helm_docs_2026}
\item \textbf{Deklarative Baseline-Applikation (Priorität 3):} Nach Registrierung der \gl{crd} \texttt{applications.argoproj.io} appliziert das Skript \texttt{application-baseline.yaml} (\texttt{baseline-workload-app}). Diese Kustomize-Applikation betreibt ein \texttt{nginx:1.30.7-alpine}-Deployment mit einem \texttt{ClusterIP}-Service im Namespace \texttt{workloads} und aktiviert \texttt{selfHeal: true}, \texttt{prune: true} sowie den Retry-Backoff aus Gleichung~\ref{eq:exponential-backoff}.\autocite{argocd_architecture_2026,kubernetes_kustomize_2026}
\end{enumerate}

Tabelle~\ref{tab:ressourcenallokation-komponenten} dokumentiert die \texttt{requests}- und \texttt{limits}-Allokation aller verwalteten Container auf dem \texttt{t3.medium}-Knoten.\autocite{burns_borg_2016,burns_kubernetes_2022}

\begin{table}[H]
    \centering
    \small
    \begin{tabularx}{\textwidth}{>{\raggedright\arraybackslash}p{3.5cm}|>{\raggedright\arraybackslash}p{1.9cm}|>{\raggedright\arraybackslash}p{2.2cm}|>{\raggedright\arraybackslash}p{2.2cm}|>{\raggedright\arraybackslash}X}
    \textbf{Komponente} & \textbf{Namespace} & \textbf{Requests (\gl{cpu}/\gl{ram})} & \textbf{Limits (\gl{cpu}/\gl{ram})} & \textbf{Exponierte Ports und Funktion} \\
    \specialrule{1.5pt}{0pt}{0pt}
    \texttt{argocd-application-\newline controller} & \texttt{argocd} & \texttt{100m} / \texttt{256Mi} & \texttt{500m} / \texttt{512Mi} & Port \texttt{8082} (Metrics); zentrales \gl{sut} und \gl{oom}-Schwellenwert \\
    \hline
    \texttt{argocd-repo-server} & \texttt{argocd} & \texttt{50m} / \texttt{128Mi} & \texttt{300m} / \texttt{384Mi} & Port \texttt{8084} (Metrics); Git-Abruf (\gl{https}) und Kustomize-Build \\
    \hline
    \texttt{argocd-server} & \texttt{argocd} & \texttt{50m} / \texttt{128Mi} & \texttt{250m} / \texttt{256Mi} & NodePort \texttt{30080} (\gl{http}/\gl{grpc}); Web-UI und \gl{api}-Endpunkt \\
    \hline
    \texttt{argocd-redis} & \texttt{argocd} & \texttt{20m} / \texttt{64Mi} & \texttt{100m} / \texttt{128Mi} & In-Memory-Cache für Manifeste und Cluster-Zustände \\
    \hline
    \texttt{prometheus-server} & \texttt{monitoring} & \texttt{50m} / \texttt{128Mi} & \texttt{300m} / \texttt{384Mi} & NodePort \texttt{30090}; hochauflösende \gl{tsdb} (Abtastung: \texttt{5s}) \\
    \hline
    \texttt{prometheus-node-\newline exporter} & \texttt{monitoring} & \texttt{10m} / \texttt{24Mi} & \texttt{50m} / \texttt{48Mi} & Host-Kernel- und \texttt{/proc/vmstat}-Metriken \\
    \hline
    \texttt{kube-state-metrics} & \texttt{monitoring} & \texttt{10m} / \texttt{32Mi} & \texttt{50m} / \texttt{64Mi} & Kubernetes-Objekt- und Pod-Zustände \\
    \hline
    \texttt{baseline-workload}\newline (\texttt{nginx:1.30.7-alpine}) & \texttt{workloads} & \texttt{20m} / \texttt{32Mi} & \texttt{100m} / \texttt{64Mi} & Port \texttt{80} (\texttt{ClusterIP}); durch ArgoCD überwachte Applikation \\
    \hline\hline
    \end{tabularx}
\caption{Ressourcenallokation (\texttt{requests} und \texttt{limits}) sowie Netzwerkschnittstellen der bereitgestellten Komponenten, eigene Darstellung.}
\label{tab:ressourcenallokation-komponenten}
\end{table}

Die Summe aller Speicher-Requests beträgt exakt \texttt{792Mi}, während die kumulierten Limits bei \texttt{1840Mi} liegen. Zusammen mit dem Grundbedarf des Ubuntu-Betriebssystems sowie der K3s-Basisschicht samt \texttt{containerd}, CoreDNS und \texttt{local-path-provisioner} (etwa 700--900~MiB) verbleibt auf der 4~GiB großen \texttt{t3.medium}-Instanz ausreichend physischer Arbeitsspeicher, um Host-seitiges Paging (\textit{Swap-Reinheit}) vollständig auszuschließen, während die ArgoCD-Control-Plane durch ihre Container-Limits (\texttt{512Mi} bzw.\ \texttt{384Mi}) weiterhin strikt ressourcenbeschränkt bleibt.\autocite{burns_borg_2016,gregg_systems_2020}

\section{Deterministische Fehlerinjektion auf Linux-Kernel-Ebene}
\label{sec:deterministische-fehlerinjektion}

Zusätzliche Kubernetes-Chaos-Operatoren wie \textit{Chaos Mesh} oder \textit{LitmusChaos} würden auf einer \texttt{t3.medium}-Instanz durch eigene Controller und Webhooks einen erheblichen Speicher- und \gl{cpu}-Overhead verursachen.\autocite{basiri_chaos_2016,koziolek_k3s_2023} Das Framework nutzt daher das schlanke Skript \texttt{inject-iptables.sh}, das Paketverluste direkt im \texttt{netfilter}-Subsystem des Linux-Kernels ohne Container-Zusatzlast erzeugt.\autocite{gregg_systems_2020,netfilter_iptables_2026}

\subsection{Paketfilterung in der Kernel-Chain FORWARD}
\label{subsec:paketfilterung-forward}

Zur Simulation einer degradierten \gl{wan}-Strecke fügt \texttt{inject-iptables.sh} für konfigurierbare Wahrscheinlichkeiten $0 < p < 1$ (z.\,B.\ $p = 0{,}80$) mittels \texttt{-m statistic --mode random --probability} eine probabilistische Bernoulli-\texttt{DROP}-Regel ein; bei $p = 1{,}00$ (vollständige Netzwerkpartition) greift ein deterministisches \texttt{-j DROP} ohne Statistikmodul.\autocite{eddy_rfc9293_2022,netfilter_iptables_2026} Da \texttt{DROP} im Gegensatz zu \texttt{REJECT} weder \texttt{RST}- noch \texttt{ICMP}-Fehlermeldungen zurücksendet, erzwingt der Linux-\gl{tcp}-Stack reale Retransmission-Timeouts (\gl{rto}).\autocite{eddy_rfc9293_2022}

Die Regel wird gezielt an Position~1 (\texttt{-I FORWARD 1}) ausschließlich in der \texttt{FORWARD}-Chain platziert:\autocite{k3s_architecture_2026,netfilter_iptables_2026}

\begin{quote}
\texttt{\small iptables -I FORWARD 1 -p tcp --dport 443 -m statistic --mode random --probability 0.80 -m comment --comment \textquotedbl{}gitops-chaos\textquotedbl{} -j DROP}
\end{quote}

Da der \texttt{argocd-repo-server} in einem isolierten Pod-Netzwerknamensraum läuft, passiert sein ausgehender Verkehr die Flannel-\gl{cni}-Bridge (\texttt{cni0}) und wird vom Kernel über \texttt{FORWARD} geroutet; Position~1 stellt sicher, dass das Verwerfen noch vor den regulären \gl{cni}-\texttt{ACCEPT}-Regeln greift.\autocite{burns_kubernetes_2022,k3s_architecture_2026} Auf eine zusätzliche Regelung der Host-\texttt{OUTPUT}-Chain wird bewusst verzichtet, um Systemdienste auf dem Host nicht mitzufiltern und den Paketverwurfzähler nicht zu verrauschen. Durch die Beschränkung auf Zielport \texttt{--dport 443} wird zudem ein streng asymmetrischer Uplink-Paketverlust auf dem ausgehenden \gl{https}-Pfad zu GitHub modelliert.

\subsection{Begrenzung des Blast Radius und fehlersichere Selbstbereinigung}
\label{subsec:blast-radius-selbstbereinigung}

Durch den Selektor \texttt{-p tcp --dport 443} in \texttt{FORWARD} bleibt der \textit{Blast Radius} strikt auf externe \gl{https}-Verbindungen der Pods begrenzt:\autocite{basiri_chaos_2016,netfilter_iptables_2026} Die \gl{ssh}-Steuerung (Port \texttt{22}), das Cluster-interne Zusammenspiel zwischen Controller, \texttt{kube-apiserver} (Port \texttt{6443}), Repo-Server und Redis sowie das lokale Prometheus-Scraping (Ports \texttt{8082}, \texttt{8084} und \texttt{30090}) bleiben vollständig ungestört.\autocite{argocd_architecture_2026,brazil_prometheus_2018}

Alle injizierten Regeln tragen den Kommentar \texttt{\textquotedbl{}gitops-chaos\textquotedbl{}}. Über \texttt{trap cleanup EXIT INT TERM} protokolliert die Bereinigungsfunktion unmittelbar vor dem Löschen den exakten Gesamtzähler verworfener Pakete (\texttt{FINAL\_DROPS} über \texttt{iptables -L FORWARD -v -n -x}) und entfernt anschließend sämtliche markierten Regeln rückstandsfrei über \texttt{iptables -S FORWARD}.\autocite{netfilter_iptables_2026}

\section{Instrumentierung, On-Node-Versuchssteuerung und Validierung der Swap-Reinheit}
\label{sec:instrumentierung-versuchssteuerung-validierung}

Um transiente Störungsfenster frei von externen \gl{wan}-Latenzen und Abtastverlusten zu erfassen, kombiniert das Framework ein 5-Sekunden-Prometheus-Scraping mit einer vollständig lokal auf der \gl{ec2}-Instanz ausgeführten Versuchssteuerung.\autocite{basiri_chaos_2016,brazil_prometheus_2018}

\subsection{Hochauflösende Metrikerfassung, Network-Namespace-Isolation und PromQL-Abfragen}
\label{subsec:metrikerfassung-promql}

In \texttt{values-prometheus.yaml} sind \texttt{scrape\_interval} und \texttt{evaluation\_interval} auf \texttt{5s} (\texttt{scrape\_timeout: 4s}) konfiguriert, um auch kurzfristige Latenz- und Ressourcenspitzen ohne \textit{Undersampling} abzubilden.\autocite{beyer_sre_2016,prometheus_docs_2026} Neben den aus kumulativen Histogrammen abgeleiteten 50.- und 95.-Perzentilen (Gleichung~\ref{eq:promql-histogram-quantile}) erfasst das Framework das exakte arithmetische Mittel (\texttt{rate(...\_sum) / rate(...\_count)}), um Bucket-Interpolationsartefakte einzuordnen:\autocite{brazil_prometheus_2018,prometheus_docs_2026}

\begin{equation}
L_{\varphi}(t) = \text{histogram\_quantile}\!\left(\varphi,\; \sum_{\text{le}} \text{rate}\!\left(\text{\texttt{argocd\_app\_reconcile\_bucket}}[1\text{m}]\right)\right)
\label{eq:promql-histogram-quantile}
\end{equation}

Dabei löst die Instrumentierung vier wesentliche Mess- und Architekturprobleme:

\begin{enumerate}
\item \textbf{Linux-Network-Namespace-Isolation (\texttt{netns}):} Der Host-\texttt{node-exporter} erfasst lediglich \texttt{/proc/net/snmp} des Host-Namespaces, während die \gl{tcp}-Retransmissions des \texttt{argocd-repo-server} ausschließlich in dessen isoliertem Pod-Namespace gezählt werden.\autocite{burns_kubernetes_2022,eddy_rfc9293_2022} Deshalb liest \texttt{node-runner.py} alle 5~Sekunden gezielt die \texttt{RetransSegs}-Zähler aus \texttt{/proc/<pid>/net/snmp} des \texttt{argocd-repo-server}-Prozesses (\texttt{pod\_tcp\_retrans\_rate}) sowie die exakten \texttt{FORWARD}-Paketverwürfe (\texttt{iptables\_drop\_rate}).\autocite{gregg_systems_2020}
\item \textbf{Keine Vorwärtsinterpolation auf Latenzen:} In inaktiven 1-Minuten-Fenstern gibt \texttt{histogram\_quantile()} \texttt{NaN} zurück. Das Framework verzichtet bewusst auf ein Vorwärtsauffüllen (\texttt{ffill}) oder Nullsetzen solcher Lücken in der \gl{csv}, um künstliche Latenzplateaus nach Störungsende zu vermeiden.
\item \textbf{Komponentenübergreifende Last-, CFS-Throttling- und Observer-Kontrolle:} Application Controller (\texttt{512Mi}) und Repo-Server (\texttt{384Mi}) werden getrennt im Namespace \texttt{argocd} samt \gl{cpu}-Throttling (\texttt{container\_cpu\_cfs\_throttled\_seconds\_total}) überwacht; zusätzlich quantifiziert der Namespace \texttt{monitoring} den Eigenbedarf der Messinfrastruktur, während \texttt{cpu\_steal\_rate} Hypervisor-Einflüsse ausschließt.\autocite{argocd_architecture_2026,gregg_systems_2020}
\item \textbf{Double-Scrape-Audit:} Zu Batch-Beginn prüft \texttt{node-runner.py} mittels \texttt{count(argocd\_git\_request\_total)}, dass Prometheus jedes Ziel genau einmal scrapt, und protokolliert sämtliche real laufenden Softwareversionen im Batch-\gl{json}.
\end{enumerate}

Tabelle~\ref{tab:prometheus-metriken-promql} fasst alle 19 \gl{promql}-Abfragen und die drei direkten Kernel-/Status-Zeitreihen zusammen.\autocite{argocd_architecture_2026,prometheus_docs_2026}

\begin{table}[H]
    \centering
    \small
    \begin{tabularx}{\textwidth}{>{\raggedright\arraybackslash}p{3.1cm}|>{\raggedright\arraybackslash}p{5.8cm}|>{\raggedright\arraybackslash}p{1.0cm}|>{\raggedright\arraybackslash}X}
    \textbf{Metrik (\gl{csv})} & \textbf{Implementierte Abfrage / Quelle} & \textbf{Einh.} & \textbf{Analytische Bedeutung} \\
    \specialrule{1.5pt}{0pt}{0pt}
    \texttt{reconcile\_duration\_\newline median\_sec} / \texttt{p95\_sec} & \texttt{\footnotesize histogram\_quantile(\{0.50,0.95\},\newline sum(rate(argocd\_app\_reconcile\_\newline bucket[1m])) by (le))} & $\text{s}$ & Median ($p_{50}$) und Tail-Latenz ($p_{95}$) der Reconciliation \\
    \hline
    \texttt{reconcile\_duration\_\newline avg\_sec} & \texttt{\footnotesize sum(rate(argocd\_app\_reconcile\_\newline sum[1m])) / sum(rate(...\_count[1m]))} & $\text{s}$ & Exaktes Mittel ohne Bucket-Interpolation \\
    \hline
    \texttt{git\_request\_\newline duration\_p95\_sec} & \texttt{\footnotesize histogram\_quantile(0.95, sum(rate(\newline argocd\_git\_request\_duration\_\newline seconds\_bucket[1m])) by (le))} & $\text{s}$ & $p_{95}$-Dauer ausgehender Git-Abrufe (\texttt{repo-server}) \\
    \hline
    \texttt{git\_request\_rate} & \texttt{\footnotesize sum(rate(argocd\_git\_request\_\newline total[1m])) or vector(0)} & $\text{s}^{-1}$ & Rate ausgehender Git-Anfragen des \texttt{repo-server} \\
    \hline
    \texttt{sync\_succeeded\_rate} / \texttt{sync\_failed\_rate} & \texttt{\footnotesize sum(rate(argocd\_app\_sync\_total\newline \{phase=\textquotedbl{}Succeeded\textquotedbl{}|\~{}\textquotedbl{}Failed|Error\textquotedbl{}\}[1m]))} & $\text{s}^{-1}$ & Erfolgs- bzw.\ Fehlerrate realer Auto-Syncs \\
    \hline
    \texttt{controller\_memory\_\newline bytes} / \texttt{cpu\_cores} & \texttt{\footnotesize working\_set\_bytes} / \texttt{\footnotesize cpu\_usage} (\texttt{\footnotesize ns=\textquotedbl{}argocd\textquotedbl{}, controller}) & $\text{B}$ /\newline \gl{vcpu} & Speicher (vs.\ \texttt{512Mi}) und \gl{cpu}-Last des Controllers \\
    \hline
    \texttt{repo\_server\_memory\_\newline bytes} / \texttt{cpu\_cores} & \texttt{\footnotesize working\_set\_bytes} / \texttt{\footnotesize cpu\_usage} (\texttt{\footnotesize ns=\textquotedbl{}argocd\textquotedbl{}, repo-server}) & $\text{B}$ /\newline \gl{vcpu} & Speicher (vs.\ \texttt{384Mi}) und \gl{cpu}-Last des Repo-Servers \\
    \hline
    \texttt{controller\_cpu\_throttled\_\newline rate} / \texttt{repo\_server\_\newline cpu\_throttled\_rate} & \texttt{\footnotesize rate(container\_cpu\_cfs\_throttled\_\newline seconds\_total[1m])} & $\text{s}/\text{s}$ & Kubernetes-CFS-\gl{cpu}-Throttling je Control-Plane-Pod \\
    \hline
    \texttt{monitoring\_memory\_\newline bytes} / \texttt{cpu\_cores} & \texttt{\footnotesize container\_memory} / \texttt{\footnotesize cpu\_usage} (\texttt{\footnotesize ns=\textquotedbl{}monitoring\textquotedbl{}}) & $\text{B}$ /\newline \gl{vcpu} & Quantifizierung des Single-Node-Observer-Overheads \\
    \hline
    \texttt{pod\_tcp\_retrans\_\newline rate} / \texttt{iptables\_\newline drop\_rate} & Direkt auf EC2: \texttt{\footnotesize /proc/<pid>/net/snmp} (\texttt{repo-server}) \& \texttt{\footnotesize iptables -L FORWARD} & $\text{s}^{-1}$ & Pod-\gl{tcp}-Retransmissions und exakte \texttt{FORWARD}-Drops \\
    \hline
    \texttt{host\_tcp\_retrans\_\newline rate} / \texttt{cpu\_steal\_rate} & \texttt{\footnotesize rate(node\_netstat\_Tcp\_RetransSegs\newline [1m])} / \texttt{\footnotesize mode=\textquotedbl{}steal\textquotedbl{}} & $\text{s}^{-1}$ /\newline \gl{vcpu} & Host-\texttt{netns}-Kontrolle und Hypervisor-Throttling-Check \\
    \hline
    \texttt{swap\_used\_bytes} / \texttt{swap\_page\_rate} & \texttt{\footnotesize SwapTotal - SwapFree} /\newline \texttt{\footnotesize rate(node\_vmstat\_pswp\{in,out\}[1m])} & $\text{B}$ /\newline $\text{s}^{-1}$ & Prometheus-Kontrolle von Swap-Nutzung und Paging \\
    \hline
    \texttt{app\_degraded\_state} & Direkt auf EC2 via \texttt{\footnotesize kubectl -o json} ($0$ oder $1$) & $[0,1]$ & Binärer Verlaufsindikator für Fehler-/\texttt{Drift}-Zustände \\
    \hline\hline
    \end{tabularx}
\caption{Im Messframework implementierte Prometheus- und On-Node-Kernel-Metriken, eigene Darstellung in Anlehnung an Argo Project (2026) und Prometheus Authors (2026).}
\label{tab:prometheus-metriken-promql}
\end{table}

\subsection{Autarke On-Node-Versuchssteuerung, Kontrollgruppe und Nachweis der Swap-Reinheit}
\label{subsec:phasengesteuerter-versuchsablauf-swap}

Um Einflüsse der Client-Netzwerkverbindung und wiederholte \gl{ssh}-Handshake-Lasten auf der \gl{ec2}-Instanz auszuschließen, überträgt \texttt{run-experiment.py} die Werkzeuge zu Beginn per \gl{scp} und startet \texttt{node-runner.py} entkoppelt (\texttt{nohup}) direkt auf dem Zielknoten.\autocite{basiri_chaos_2016,gregg_systems_2020} Alle Zeitstempel und Kernel-Auslesungen erfolgen damit auf einer gemeinsamen Systemuhr. Neben verschiedenen Paketverluststufen (Parameter-Sweep bis $p = 1{,}00$) unterstützt \texttt{node-runner.py} echte Kontrollläufe (\texttt{--probability 0.0}) mit identischem Trigger- und Drift-Protokoll. Sequenzielle Wiederholungsläufe werden durch einen 60-sekündigen Cooldown getrennt, damit gleitende 1-Minuten-\gl{promql}-Fenster vor dem nächsten Lauf abklingen.\autocite{brazil_prometheus_2018} Jeder Lauf gliedert sich in fünf Phasen:

\begin{enumerate}
\item \textbf{Phase 0 -- Vorbereitung:} \texttt{run-experiment.py} startet \texttt{node-runner.py} lokal auf der \gl{ec2}-Instanz. Dieser bereinigt etwaige Altregeln in \texttt{iptables}, prüft die lokale Prometheus-Erreichbarkeit, verifiziert über \texttt{kubectl} den stabilen Ausgangszustand (\texttt{Synced}, \texttt{Healthy}, kein Drift-Label) und protokolliert die Systemumgebung.
\item \textbf{Phase 1 -- Steady-State-Baseline ($30\,\text{s}$):} Unmittelbar vor Beginn der Baseline liest \texttt{node-runner.py} die Zähler $\text{pswpin}(t_{\text{start}})$ und $\text{pswpout}(t_{\text{start}})$ direkt aus \texttt{/proc/vmstat} ein und zeichnet 30~Sekunden lang im 5-Sekunden-Takt den ungestörten Betriebszustand auf.\autocite{basiri_chaos_2016,kernel_vmstat_2026}
\item \textbf{Phase 2 -- Fehlerinjektion, Cache-Invalidierung und Drift ($120\,\text{s}$):} Zum Zeitpunkt $t_{\text{fault\_start}}$ aktiviert \texttt{node-runner.py} in Chaos-Läufen \texttt{inject-iptables.sh}. Bei $t = +3\,\text{s}$ (sowie erneut alle 30~Sekunden während des Störungsfensters) erzwingt er per \texttt{argocd.argoproj.io/refresh=hard} eine Cache-Invalidierung und injiziert erst 1~Sekunde danach ($t = +4\,\text{s}$) den Label-Drift \texttt{experiment=drifted} auf dem Deployment, damit ein gecachtes Manifest keinen sofortigen Self-Heal über Port \texttt{6443} auslöst.\autocite{argocd_architecture_2026} Im 5-Sekunden-Takt protokolliert das Skript den Eintritt eines Netzfehlers ($t_{\text{fault\_detected}}$ bei \texttt{ComparisonError} bzw.\ \texttt{Unknown}), ignoriert dabei etwaige transiente \texttt{kubectl}-Fehler und erfasst über einen Waiter-Thread das Störungsende $t_{\text{fault\_end}}$.
\item \textbf{Phase 3 -- Konvergenz, Nachlaufzeit ($90\,\text{s}$) und Swap-Reinheit:} Nach $t_{\text{fault\_end}}$ überwacht das Skript (Timeout: 240~Sekunden) den Abschluss der Konvergenz ($t_{\text{recovered}}$), wobei neben \texttt{Synced} und \texttt{Healthy} explizit die Bereinigung des Drift-Labels geprüft wird. Eine anschließende 90-sekündige Stabilisierungsphase stellt sicher, dass alle 1-Minuten-\gl{promql}-Fenster in der \texttt{post\_recovery}-Phase vollständig abklingen. Danach werden $T_{\text{dwell}}$ (Gleichung~\ref{eq:dwell-time}), $\text{\gl{mttr}}_{\text{post-fault}}$ (Gleichung~\ref{eq:mttr}) und die formale Bedingung der Swap-Reinheit geprüft:\autocite{beyer_sre_2016,gregg_systems_2020,kernel_vmstat_2026}
\begin{equation}
\begin{aligned}
\Delta\text{pswpin} &= \text{pswpin}(t_{\text{end}}) - \text{pswpin}(t_{\text{start}}) \stackrel{!}{=} 0 \\
\Delta\text{pswpout} &= \text{pswpout}(t_{\text{end}}) - \text{pswpout}(t_{\text{start}}) \stackrel{!}{=} 0
\end{aligned}
\label{eq:swap-purity-condition}
\end{equation}
Abschließend exportiert \texttt{node-runner.py} alle Zeitreihen inklusive \texttt{experiment\_phase} (mit echten \texttt{NaN}-Lücken bei Inaktivität) in eine \gl{csv}-Datei sowie phasengetrennte und aggregierte Statistiken in eine \gl{json}-Zusammenfassung.
\item \textbf{Phase 4 -- Artefakt-Transfer und Visualisierung:} Nach dem \texttt{DONE}-Signal lädt \texttt{run-experiment.py} alle Ergebnisdateien per \gl{scp} herunter. Das Skript \texttt{analysis/plot-results.py} generiert daraus Zeitreihen- und Ressourcenverläufe (mit schattierten Fehler-/Erholungsphasen und Observer-Overhead) sowie bei Mehrfachläufen bzw.\ Parameter-Sweeps Cross-Run-Boxplots über wiederholte Messläufe.\autocite{brazil_prometheus_2018}
\end{enumerate}


    % Abkürzungsverzeichnis
    \newpage
    \phantomsection
    \label{sec:Abkuerzungsverzeichnis}
    \addcontentsline{toc}{chapter}{Abkürzungsverzeichnis}
    \printnoidxglossary[type=\acronymtype,title={Abkürzungsverzeichnis}, nonumberlist]
    \newpage
    
    % Abbildungsverzeichnis
    \phantomsection
    \addcontentsline{toc}{chapter}{Abbildungsverzeichnis}
    \markboth{Abbildungsverzeichnis}{}
    \listoffigures
    \newpage
    
    % Tabellenverzeichnis
    \phantomsection
    \addcontentsline{toc}{chapter}{Tabellenverzeichnis}
    \markboth{Tabellenverzeichnis}{}
    \listoftables
    \newpage
    
    % Promptverzeichnis (für KI-Prompts)
    %\phantomsection
    %\addcontentsline{toc}{chapter}{Promptverzeichnis}
    %\markboth{Promptverzeichnis}{} % Manuelles Setzen der Kopfzeile
    %\listofprompt
    %\newpage

    % Literaturverzeichnis
    \newpage
    \phantomsection
    \addcontentsline{toc}{chapter}{Literaturverzeichnis}
    \label{chap:Literaturverzeichnis}
    \markboth{Literaturverzeichnis}{} % Manuelles Setzen der Kopfzeile
    \printbibliography[title={Literaturverzeichnis}]
    
    % Anhangsverzeichnis
    %\newpage
    %\phantomsection
    %\addcontentsline{toc}{chapter}{Verzeichnis der Anhänge}
    %\listofappendix
    %\newpage
    
    % Anhänge
    %\newpage
    %\appendix
    % Spezielle Formatierung für Anhänge
    %\titleformat{\chapter}
    %  {\fontsize{11}{16.5}\bfseries\raggedright}
    %  {Anhang \thechapter:}{0.5em}{}
    % Spezielle Formatierung für Anhänge in der Kopfzeile
    %\renewcommand{\chaptermark}[1]{\markboth{Anhang \thechapter: #1}{}}
    % Spezielle Formatierung für Anhänge im Inhaltsverzeichnis
    %\addtocontents{toc}{\protect\renewcommand{\protect\cftchappresnum}{Anhang }}
    %\addtocontents{toc}{\protect\renewcommand{\protect\cftchapaftersnum}{: }}
    %\addtocontents{toc}{\protect\setlength{\protect\cftchapnumwidth}{5.5em}}
    %% Include all appendix files
\chapter{Praktische Hinweise zum Gendern}
\appendixentry{Praktische Hinweise zum Gendern}
\label{app:gendern}

Eine häufig verwendete Möglichkeit ist das substantivierte Partizip Präsens. Es zeigt im Plural kein Genus:
\begin{itemize}
\item die Studierenden,
\item die Lehrenden.
\end{itemize}

Bitte beachten Sie, dass das substantivierte Partizip Präsens im Singular ein Genus zeigt:
\begin{itemize}
\item der Studierende ist eindeutig maskulin,
\item die Lehrende ist eindeutig feminin.
\end{itemize}

Abstrakte Rollenbezeichnungen und Begriffe für Kollektive, Institutionen und Positionen sind eine weitere Möglichkeit für sexusindifferente Formulierungen:
\begin{itemize}
\item die Vertretung,
\item die Studiengangsleitung,
\item das Kollegium,
\item das Publikum,
\item die Lehrkräfte,
\item das Team,
\item das Präsidium.
\end{itemize}

Die Pronomina alle und niemand zeigen ebenfalls kein Genus:
\begin{itemize}
\item Alle können am Kurs teilnehmen.
\item Niemand muss sich wegen der Prüfung Sorgen machen.
\end{itemize}

Gender-Zeichen, meist ein Doppelpunkt oder Sternchen (Asterisk), stehen zwischen dem Wortstamm der männlichen Form und der weiblichen Endung:
\begin{itemize}
\item Bewerber:in
\item Teilnehmer:innen
\item Student:in
\item Expert:innen
\end{itemize}

Die Beidnennung ist eine weitere Möglichkeit, also die Nennung von maskulinen und femininen Bezeichnungen:
\begin{itemize}
\item Studentinnen und Studenten,
\item Dozentinnen und Dozenten,
\item Prüferin bzw. Prüfer,
\item Expertin bzw. Experte.
\end{itemize}
\chapter{Eigenständigkeitserklärung}
\appendixentry{Eigenständigkeitserklärung}
\label{app:eigenstaendigkeitserklaerung}

Ihre Eigenständigkeitserklärung hat folgenden Wortlaut (die englische Version verwenden Sie nur, wenn Sie Ihre Arbeit auf Englisch verfassen). Bitte denken Sie daran, sie eigenhändig zu unterschreiben.

\subsection*{Eigenständigkeitserklärung}

Hiermit bestätige ich, dass ich die vorliegende Arbeit selbstständig verfasst und keine anderen Publikationen, Vorlagen und Hilfsmittel (z.B. künstliche Intelligenz) als die angegebenen benutzt habe. Alle Teile meiner Arbeit, die wortwörtlich oder dem Sinn nach anderen Werken entnommen sind, wurden unter Angabe der Quelle kenntlich gemacht. Gleiches gilt für von mir verwendete Internetquellen. Ich versichere, dass ich diese Arbeit oder nicht zitierte Teile daraus vorher nicht in einem anderen Prüfungsverfahren eingereicht habe. Mir ist bekannt, dass meine Arbeit zum Zwecke eines Plagiatsabgleichs mittels einer Plagiatserkennungssoftware auf eine ungekennzeichnete Übernahme von fremdem geistigen Eigentum sowie auf die Nutzung von künstlicher Intelligenz zur Texterstellung überprüft werden kann. Ich versichere, dass die elektronische Form meiner Arbeit mit der gedruckten Version identisch ist.

Datum, Unterschrift

\subsection*{Declaration of Originality}

I hereby confirm that I have independently written this work and have not used any publications, templates, or aids (e.g. artificial intelligence) other than those I have indicated. All parts of my work which have been taken literally or correspondingly from other publications have been duly acknowledged. This also applies to Internet sources. I confirm that I have not previously submitted this work or any unquoted parts thereof in any other examination procedure. I am aware that my work may be checked for plagiarism by means of plagiarism recognition software, as well as for the use of artificial intelligence for text creation, in order to verify the integrity of its written content. I also confirm that the electronic form is identical to the printed version.

Date, signature
\chapter{Ergänzende Hinweise für Abschlussarbeiten in den Studiengängen Wirtschaftspsychologie und Mediendesign}
\appendixentry{Ergänzende Hinweise für Abschlussarbeiten für Abschlussarbeiten in den Studiengängen Wirtschaftspsychologie und Mediendesign}
\markboth{Anhang C: Ergänzende Hinweise}{}
\label{app:wirtschaftspsychologie-mediendesign}

In diesen beiden Studiengängen können die Abschlussarbeit unterschiedlich ausgerichtet werden. Möglich sind:
\begin{enumerate}
\item theoretisch-konzeptionelle Abschlussarbeiten,
\item empirische Arbeiten (in der Wirtschaftspsychologie verpflichtend): qualitative oder quantitative Untersuchungen,
\item konzeptionell-gestalterische Arbeiten (nur im Studiengang Mediendesign).
\end{enumerate}

Eine eindeutige Zuordnung zu einem dieser Schwerpunkte ist nicht immer möglich.

In jedem Fall ist eine umfassende und systematische Darstellung des Stands der Forschung anhand der aktuellen wissenschaftlichen Literatur (Diskurs) erforderlich. Auf diesen Status Quo bezieht sich Ihre weitere Bearbeitung, bzw. Sie wenden ihn dafür an.

Alle Arbeiten können in Kooperation mit Unternehmen geschrieben werden. Bitte stimmen Sie sich dazu mit Ihren Prüfern ab.

\section*{1. Theoretisch-konzeptionelle Arbeiten}

Wenn Sie eine theoretisch-konzeptionelle Arbeit vorlegen, analysieren Sie darin ein Thema, ohne dass Sie hierzu eigene Berechnungen durchführen.

Ihr Ausgangspunkt ist eine möglichst vollständige Erfassung der relevanten Literatur und die systematische Aufbereitung der gewonnenen Informationen. Wie Sie dieses Material strukturieren und aufbereiten (Tiefe, Differenzierung, Präzision, Prägnanz, Gewichtung), bestimmt wesentlich die Qualität Ihrer Arbeit.

Typischerweise wird von Ihnen erwartet, dass Sie die wesentliche Literatur durchdringen und sich mit ihr kritisch auseinandersetzen. Dazu zählen insbesondere Definitionen (die Sie evtl. selbst entwickeln), Klassifikationsschemata oder Typologien. Zudem wird erwartet, dass sie Begründungszusammenhänge strukturierte darstellen.

Ein zentrales Ziel der theoretisch-konzeptionellen Arbeit ist die Verbesserung, Vereinfachung und vor allem Systematisierung der Darstellung bereits bekannter, aber bisher nicht in einer einzigen Arbeit zusammengefasster Fakten oder Argumente.

Den Titel Ihrer theoretisch-konzeptionellen Arbeit können Sie beispielsweise nach diesem Schema formulieren:
\begin{itemize}
\item Entwicklung einer strategisch-konzeptionellen Empfehlung
\item Formulierung von Handlungsempfehlungen für bestimmte Fälle
\item Eingrenzung von potenziellen Entwicklungsszenarien
\item Konzeption von bestimmten Maßnahmenempfehlungen für besondere Zwecke (beispielsweise aus dem konkreten Arbeitsumfeld).
\end{itemize}

Ihre Strukturierungsleistung soll zu einem Ziel führen, das bereits im Titel anklingt, z.B.: \glqq Strategische Markt- und Wettbewerbsanalyse der Videospielbranche - unter besonderer Berücksichtigung der Neuprodukteinführung eines Konsolenherstellers\grqq{}.

\section*{2. Empirische Arbeiten}

Bei Ihrer empirischen Arbeit kann es sich entweder um eine qualitative oder quantitative Untersuchung handeln.

Qualitative Untersuchungen werden durchgeführt, wenn ein wesentliches Verständnis für eine bestimmte Forschungsfrage erlangt werden soll. Ziel der qualitativen Forschung ist das Erkennen, Beschreiben und Verstehen von Zusammenhängen. Im Vordergrund stehen die vollständige Erfassung und Interpretation aller problemrelevanten Aspekte hinsichtlich des Themas der Abschlussarbeit. Hierbei bedient man sich offener, nicht standardisierter Erhebungsverfahren (z.B. Interviews, Expertengespräche, Gruppendiskussionen, qualitative Beobachtung, qualitative Experimente), deskriptiver Aufbereitungsverfahren (z.B. Gesprächsprotokoll, Transkription) und interpretativer Auswertungsverfahren (z.B. qualitative Inhaltsanalyse, Cognitive Mapping, objektive Hermeneutik). In der Regel kommen im Rahmen der qualitativen Untersuchung Methoden zum Einsatz, die sich auf kleine Fallzahlen beschränken, keine statistischen Analysen (z.B. Signifikanztests) implizieren, relative weiche Daten produzieren und ihre Erkenntnisse auf einem verhältnismäßig niedrigen Abstraktionsniveau mittels subjektiver Interpretation gewinnen. Die gewonnenen Ergebnisse sind zwar nicht repräsentativ, sie dienen aber zur ausführlichen Stoffsammlung, um ihrerseits wieder Hypothesen quantifizieren zu können.

Quantitative Untersuchungen haben die Messung bestimmter Sachverhalte bzw. die Entdeckung von Gesetzmäßigkeiten zum Gegenstand. Es handelt sich hierbei um einen Ansatz, der theoriegeleitet ist und der sich standardisierter Erhebungsmethoden (i.d.R. schriftliche Fragebögen) bedient. Ziel dieser Standardisierung ist es, die Antworten einer Vielzahl von Befragten unmittelbar vergleichen zu können. Der Vorteil quantitativer Untersuchungen liegt darin, dass sich die Messergebnisse mit statistischen Methoden (z.B. Korrelations- und Regressionsanalysen, Varianzanalysen, Faktoren- und Clusteranalysen, Metaanalysen, Diskriminanzanalysen, MDS) unter Nutzung statistischer Kennzahlen (Mittelwert, Median, Standardabweichung, Varianz etc.) verdichten und weiterverarbeiten lassen. Die Auswertung und Analyse der Daten, die im Rahmen der quantitativen Forschung generiert werden, erfolgt mit Hilfe von statistischen Programmen (z.B. SPSS, AMOS, PLS, R). Ziel ist es, Zusammenhänge zu erkennen und daraus allgemein gültige Aussagen abzuleiten. Rückschlüsse auf die tatsächlichen Verhältnisse in der Grundgesamtheit sind möglich.

In der Wirtschaftspsychologie (\gl{bsc}) werden ausschließlich empirische Abschlussarbeiten durchgeführt, wobei quantitativen Ansätzen der Vorzug zu geben ist.

\section*{3. Konzeptionell-gestalterische Arbeiten}

Diese Arbeit können Sie nur im Studiengang Mediendesign vorlegen. Hierbei handelt es sich um eine innovative Lösung für eine komplexe gestalterische Aufgabe.

Ihr Ausgangspunkt ist auch hier die Durchdringung relevanter wissenschaftlicher Grundlagen (insbesondere Theorien und Beispiele aus der Designgeschichte).

Daran schließen sich evtl. strategische sowie Ihre konzeptionellen und methodischen Darstellungen an, in denen Sie den Kontext der Aufgabe und Ihren Designprozess nachvollziehbar machen.

Sowohl die wissenschaftlichen als auch die methodischen und konzeptionellen Aussagen führen zu Ihrer praktischen Umsetzung im Sinn einer Lösungsentwicklung mit gestalterischen und technischen Mitteln.
\chapter{Gestalterische Hinweise für die Studiengänge Mediendesign (B.A.) und User Experience Design (M.A.)}
\appendixentry{Gestalterische Hinweise für die Studiengänge Mediendesign (B.A.) und User Experience Design (M.A.)}
\markboth{Anhang D: Gestalterische Hinweise}{}
\label{app:gestalterische-hinweise}

\section*{Schmutztitel und Haupttitel}

Auf dem Schmutztitel und dem Haupttitel stehen die gleichen Inhalte wie bei allen anderen Studiengängen (vgl. Abschnitt 6.6 in diesem Leitfaden). Ansonsten sind Sie in Ihrer Gestaltung der Titel frei. Es ist auch ausdrücklich erwünscht, dass Sie von dieser Freiheit Gebrauch machen. Falls Sie sich dafür entscheiden, Ihre Arbeit nicht zu gestalten, verwenden Sie einen braunen (Mediendesign) bzw. weißen (User Experience Design) Einband. Bitte beachten Sie, dass Sie keinen schlichten andersfarbigen Einband verwenden, weil diese bereits anderen Studiengängen vorbehalten sind.

\section*{Schriften}

Beachten Sie bei der Wahl Ihrer Schrift die Angemessenheit. Neben der akademischen Ausstrahlung ist es wichtig, dass Sie sich später als Bewerber mit Ihrer Arbeit gegen Konkurrenten durchsetzten wollen. Alle Schriften für sogenannte Mengentexte sind aufgrund ihrer Lesefreundlichkeit geeignet, beispielsweise die Thesis, Info Text oder Corporate oder Schriften nach dem dynamischen Formprinzip, wie z.B. Garamond, Minion, Weidemann, Lexicon, Optima, Officina, Frutiger, Today oder Meta.

Beachten Sie bitte auch, dass die Prüfer Ihre Arbeit innerhalb kurzer Zeit lesen werden. Die Lektüre darf auch Freude bereiten, deshalb ist die Lesbarkeit bei der Auswahl ein wichtiges Entscheidungskriterium. Eine schmal laufende Schrift (condensed) ist also nicht geeignet. Serifenschriften bieten in der Regel Vorteile, sind aber kein Muss.

\section*{Raster und Layout}

Die Bachelor-Thesis soll angemessene Seitenränder haben und nur eine Textspalte, sowohl einen toten als auch einen lebenden Kolumnentitel integrieren sowie zwingend über eine Paginierung verfügen. Abbildungen sollt angemessen groß sein. Komplexe Schaubilder oder die Herleitung Ihrer Entwürfe dürfen, sofern notwendig, seitenfüllend sein. Für Infografiken entwickeln Sie idealerweise einen durchgängigen einheitlichen Stil.

\section*{Feintypografie}

Beachten Sie bitte alle feintypografischen Regeln. Halten Sie bitte sauber definierte Abstände vor und nach Überschriften oder Absätzen usw. ein. Ein Ausgleich des Flatter- oder Blocksatzes nach Fertigstellen des Textes ist Pflicht. Tipp: Beginnen Sie damit von hinten.

    % Eigenständigkeitserklärung
    \newpage
    \markboth{Eigenständigkeitserklärung}{}
    \chapter*{Eigenständigkeitserklärung}
\addcontentsline{toc}{section}{Eigenständigkeitserklärung}
Hiermit bestätige ich, dass ich die vorliegende Arbeit selbstständig verfasst und keine 
anderen Publikationen, Vorlagen und Hilfsmittel (z.B. künstliche Intelligenz) als die 
angegebenen benutzt habe. Alle Teile meiner Arbeit, die wortwörtlich oder dem Sinn nach 
anderen Werken entnommen sind, wurden unter Angabe der Quelle kenntlich gemacht. 
Gleiches gilt für von mir verwendete Internetquellen. Ich versichere, dass ich diese Arbeit 
oder nicht zitierte Teile daraus vorher nicht in einem anderen Prüfungsverfahren eingereicht 
habe. Mir ist bekannt, dass meine Arbeit zum Zwecke eines Plagiatsabgleichs mittels einer 
Plagiatserkennungssoftware auf eine ungekennzeichnete Übernahme von fremdem geistigen 
Eigentum sowie auf die Nutzung von künstlicher Intelligenz zur Texterstellung überprüft 
werden kann. Ich versichere, dass die elektronische Form meiner Arbeit mit der gedruckten 
Version identisch ist.


% \section*{Declaration of Originality}
% I hereby confirm that I have independently written this work and have not used any 
% publications, templates, or aids (e.g. artificial intelligence) other than those I have 
% indicated. All parts of my work which have been taken literally or correspondingly from other
% publications have been duly acknowledged. This also applies to Internet sources. I confirm
% that I have not previously submitted this work or any unquoted parts thereof in any other 
% examination procedure. I am aware that my work may be checked for plagiarism by means 
% of plagiarism recognition software, as well as for the use of artificial intelligence for text
% creation, in order to verify the integrity of its written content. I also confirm that the 
% electronic form is identical to the printed version. \\

\vspace{2cm}

\begin{center}
\begin{tabular}{p{0.45\textwidth}p{0.45\textwidth}}
    [Ort], \today & \\
    \hline
    % Place, Date & Signature \\
    Ort, Datum & Unterschrift \\
\end{tabular}
\end{center}
    
    % Lebenslauf (nur für Abschlussarbeiten)
    \newpage
    \markboth{Lebenslauf}{}
    \chapter*{Lebenslauf}
\addcontentsline{toc}{section}{Lebenslauf}
\label{sec:lebenslauf}

\begin{tabular}{>{\raggedright\arraybackslash}p{0.34\textwidth}>{\raggedright\arraybackslash}p{0.61\textwidth}}
\textbf{Persönliche Daten} & \\
\\
Name: & [Vorname Nachname] \\
Geburtsdatum: & [TT.MM.JJJJ] \\
Geburtsort: & [Geburtsort] \\
Staatsangehörigkeit: & [Staatsangehörigkeit] \\
\\
\textbf{Ausbildung} & \\
\\
MM.JJJJ -- heute & Studium Informatik (B.Sc.) an der Rheinischen Hochschule Köln \\
MM.JJJJ -- MM.JJJJ & Vorherige Ausbildung/Studium \\
MM.JJJJ -- MM.JJJJ & Schulbildung \\
\\
\textbf{Berufliche Erfahrung} & \\
\\
MM.JJJJ -- heute & Aktuelle Tätigkeit \\
MM.JJJJ -- MM.JJJJ & Vorherige Tätigkeit \\
MM.JJJJ -- MM.JJJJ & Praktikum/Werkstudententätigkeit \\
\\
\textbf{Kenntnisse und Fähigkeiten} & \\
\\
Sprachen: & Sprache 1 (Niveau), Sprache 2 (Niveau), ... \\
EDV-Kenntnisse: & Software 1, Software 2, ... \\
Sonstige Qualifikationen: & Qualifikation 1, Qualifikation 2, ... \\
\\
\end{tabular}

\vspace{1cm}

\begin{center}
Ort, Datum
\end{center}


\end{onehalfspace}
\end{document}
